\documentclass[11pt]{article}
\usepackage[margin=1in]{geometry}
\usepackage{amsmath,amssymb,amsthm}
\usepackage{mathtools}
\usepackage{bm}
\usepackage{array}
\usepackage{booktabs}
\usepackage{algorithm}
\usepackage{algpseudocode}
\usepackage{titlesec}
\usepackage[colorlinks=true,linkcolor=blue,citecolor=blue,urlcolor=blue]{hyperref}
\usepackage{cite}
\usepackage{enumerate}
\usepackage{xcolor}
\usepackage{microtype}

%% Numbered \section headers flow continuously by default.
%% Starred \section* uses a separate format.
\titleformat{\section}{\normalfont\Large\bfseries}{\thesection}{1em}{}
\titleformat{name=\section,numberless}{\normalfont\Large\bfseries}{}{0pt}{}

%% ── Theorem environments ──────────────────────────────────────────────────────
\newtheorem{theorem}{Theorem}[section]
\newtheorem{proposition}[theorem]{Proposition}
\newtheorem{lemma}[theorem]{Lemma}
\newtheorem{corollary}[theorem]{Corollary}
\theoremstyle{definition}
\newtheorem{definition}[theorem]{Definition}
\newtheorem{example}[theorem]{Example}
\theoremstyle{remark}
\newtheorem{remark}[theorem]{Remark}

%% ── Convenience macros ────────────────────────────────────────────────────────
\newcommand{\R}{\mathbb{R}}
\newcommand{\Lsup}{L_{\mathrm{sup}}}
\newcommand{\LSH}{L_{\mathrm{SH}}}
\newcommand{\Ksig}{\Sigma}
\newcommand{\Ctil}{\widetilde{C}}
\newcommand{\Gom}{G_\omega}
\newcommand{\Pp}{P_{p,r}}
\newcommand{\Up}{U_{p,r}}
\newcommand{\vhat}{\hat{v}}
\newcommand{\Rp}{R_p}
\newcommand{\calK}{\mathcal{K}}
\newcommand{\calC}{\mathcal{C}}
\newcommand{\calG}{\mathcal{G}}
\DeclareMathOperator{\diag}{diag}
\DeclareMathOperator{\im}{im}
\DeclareMathOperator{\spn}{span}
\DeclareMathOperator{\sgn}{sgn}
\DeclareMathOperator{\simop}{sim}
\DeclareMathOperator{\tri}{tri}

%% ─────────────────────────────────────────────────────────────────────────────

\title{The Supra-Hodge Laplacian: A Unified Spectral Operator\\
       for Multi-View Higher-Order Relational Systems}

\author{Mason Hughes\\[4pt]
        Independent Researcher\\
        Missouri, USA\\
        \texttt{masonhughesmwh@gmail.com}}

\date{}

\begin{document}
\maketitle

%% ── Abstract ─────────────────────────────────────────────────────────────────
\begin{abstract}
We introduce the \emph{Supra-Hodge Laplacian}, a new spectral operator that
unifies two previously separate generalizations of the classical graph
Laplacian: the supra-Laplacian, which encodes cross-layer coupling in
multilayer graph systems, and the combinatorial Hodge Laplacian, which encodes
higher-order topological structure within a single simplicial complex. Fixing a
family of complexes on a shared vertex set (one view per layer), an aligned
reference set of $p$-simplices, and nonnegative scalar couplings between layers,
the order-$p$ Supra-Hodge Laplacian $L_p$ is the block operator on stacked
ambient $p$-chains that augments embedded layerwise Hodge Laplacians
$H_p^{(i)}$ with a Kronecker-type coupling $L_{\mathrm{inter}}\otimes I_{N_p}$
from the inter-layer graph. For a stacked signal, its quadratic energy splits
into within-layer topological smoothness and cross-layer disagreement. The
central formal result is the
\emph{Consensus-Harmonic Nullspace Theorem}: $\ker(L_p)$ is exactly the set of
stacked signals that are simultaneously harmonic in every layer and mutually
consensual across every connected component of the inter-layer coupling graph.

The construction is provably consistent: in the single-layer, order-$0$,
no-higher-order limit it recovers the standard graph Laplacian; at $p=0$ it
reduces to the classical supra-Laplacian; and at $k=1$ it reduces to the
single-layer Hodge Laplacian. Spectral analysis yields a
generalized Fiedler theory for coupled higher-order systems, with low-energy
eigenmodes identifying coherent multi-view harmonic structure and
conflict-frontier modes identifying topological inconsistency across views. We
develop a low-energy projector $P_{p,r}$ that achieves principled spectral
compression of multi-view chain signals. An application to multi-view
large-language-model reasoning formalizes the operator as a diagnostic and
compressive device over structured reasoning traces.
Proof-of-concept experiments on synthetic relational benchmarks support the
theoretical predictions and show advantages over pairwise and single-view
spectral baselines in this controlled setting.
\end{abstract}

\noindent\textbf{Keywords:} Hodge Laplacian, supra-Laplacian, simplicial complexes,
multilayer graphs, spectral methods, higher-order topology, multi-view learning,
LLM reasoning, spectral compression.

\clearpage
\tableofcontents
\clearpage

\section*{Notation Summary}
\addcontentsline{toc}{section}{Notation Summary}

\begin{center}
\begin{tabular}{ll}
\toprule
Symbol & Meaning \\
\midrule
$K^{(i)}$ & Simplicial complex for layer $i$ \\
$H_p^{(i)}$ & Layerwise Hodge Laplacian on the aligned ambient space $\Ctil_p$ \\
$N_p$ & Ambient $p$-chain dimension, $N_p=|\Sigma_p|$ \\
$\Ctil_p$ & Common aligned ambient $p$-chain space, $\Ctil_p\cong\R^{N_p}$ \\
$L_{\mathrm{inter}}$ & Inter-layer graph Laplacian \\
$L_p$ & Supra-Hodge Laplacian \\
$v^{(i)}$ & $p$-chain signal in layer $i$ \\
$\Sigma_p$ & Global simplex index set \\
\bottomrule
\end{tabular}
\end{center}

%% ══════════════════════════════════════════════════════════════════════════════
\clearpage
\section{Introduction}
\label{sec:intro}
%% ══════════════════════════════════════════════════════════════════════════════

\subsection{Motivation: Multi-View and Higher-Order Relational Systems}
\label{subsec:motivation}

Modern relational systems are neither pairwise nor single-view. A
protein-interaction network simultaneously encodes physical binding, regulatory
co-expression, and genetic interaction, each of which can form higher-order
complexes involving three or more proteins jointly. A multi-hop reasoning agent
maintains semantic, evidential, causal, and temporal layers over the same set
of working entities, and within each layer certain groups of three entities form
genuine triplets that cannot be decomposed into their pairwise parts. A
knowledge graph for scientific planning contains both citation layers and
experimental-dependency layers, each populated with cliques and triangles that
carry structurally distinct information.

Two powerful spectral frameworks have been developed to address these
complexities individually, but neither addresses both at once. The
\emph{supra-Laplacian} framework~\cite{gomez2013,dedomenico2013} couples
multiple graph layers over a shared vertex set via a block operator, enabling
multi-view spectral analysis of pairwise relations. The \emph{combinatorial
Hodge Laplacian}~\cite{horak2013,schaub2020,lim2020} generalizes the graph
Laplacian to simplicial complexes, enabling spectral analysis of higher-order
interactions within a single relational view. Neither framework, in its usual form, is aimed at the full combination:
multiple views, each possessing higher-order structure, coupled across views.

We construct the \emph{Supra-Hodge Laplacian}, a block operator that is
simultaneously a supra-Laplacian across views and a Hodge Laplacian within each
view. The primary contribution of this paper is
mathematical; the reasoning application in Section~\ref{sec:application} serves
to illustrate the operator's relevance to structured multi-view inference
systems.

\subsection{Why Pairwise Graph Methods Are Insufficient}
\label{subsec:insufficient}

Consider three agents forming a stable epistemic coalition whose stability is a
genuinely three-body phenomenon: remove any single edge and the coalition
dissolves, yet the pairwise adjacency between any pair is nonzero. A graph
Laplacian encodes only the pairwise structure; it cannot represent the
qualitatively different status of the filled triangle versus the hollow cycle
of three edges. This is not a quirk: it is the defining limitation of all
pairwise methods.

Within a single Hodge framework, this limitation is overcome for one relational
view. But when $k$ views exist---each with its own triangles, its own harmonic
cycles, its own curl-free and gradient-free signals---a single complex conflates
their topologies. Analyzing each view in isolation misses cross-view
interactions; analyzing all views in a single complex erases the boundaries
between views.

The supra-Laplacian handles cross-view interactions, but only at the level of
vertices and pairwise edges. It cannot measure higher-order inconsistency
within a view, nor can it detect the collapse of a harmonic cycle in one view
due to pressure from a topologically simpler view coupled to it.

\subsection{What This Paper Contributes}
\label{subsec:contributions}

This paper makes the following contributions.
\begin{enumerate}
\item \textbf{Operator construction.} We define the order-$p$ Supra-Hodge
      Laplacian $L_p$ in Section~\ref{sec:defn} as a block operator on a common
      aligned $p$-chain space, combining embedded layerwise Hodge Laplacians
      $H_p^{(i)}$ with an inter-layer coupling term induced by the layer graph
      Laplacian $L_{\mathrm{inter}}$.
\item \textbf{Main theorem.} We prove the Consensus-Harmonic Nullspace
      Theorem (Theorem~\ref{thm:main}), giving a complete algebraic
      characterization of $\ker(L_p)$.
\item \textbf{Reduction propositions.} We prove that $L_p$ reduces to the
      classical supra-Laplacian at $p=0$ and to the single-layer Hodge
      Laplacian at $k=1$ (Propositions~\ref{prop:k1} and~\ref{prop:p0}).
\item \textbf{Spectral theory.} We develop a generalized Fiedler theory and
      formalize the low-energy projector $P_{p,r}$ for spectral compression
      (Section~\ref{sec:spectral}).
\item \textbf{Application.} We formalize the mapping of the framework to
      multi-view LLM reasoning as an illustrative application domain
      (Section~\ref{sec:application}).
\item \textbf{Experiments.} We present synthetic proof-of-concept experiments
      that test the operator's predicted behavior in controlled settings
      (Section~\ref{sec:experiments}).
\end{enumerate}

\subsection{Why the Paper Is Self-Contained}
\label{subsec:selfcontained}

Sections~\ref{sec:linalg}--\ref{sec:hodge} build all necessary prerequisites
from scratch: linear algebra, graph Laplacians, supra-Laplacians, simplicial
complexes, chain spaces, boundary operators, and single-layer Hodge Laplacians.
No algebraic topology background beyond what is defined in these sections is
assumed. A reader with undergraduate linear algebra and basic familiarity with
graphs should be able to follow every proof by reading sequentially.

\subsection{Roadmap for the Reader}

Sections~\ref{sec:linalg}--\ref{sec:hodge} lay the mathematical foundations.
Section~\ref{sec:coupling} constructs the alignment and coupling infrastructure;
Section~\ref{sec:defn} then defines the Supra-Hodge Laplacian $L_p$ itself.
Section~\ref{sec:properties} establishes fundamental properties and the main
theorem. Section~\ref{sec:spectral} develops spectral interpretation.
Section~\ref{sec:dynamics} treats dynamics and perturbations.
Section~\ref{sec:application} maps the theory to LLM reasoning.
Section~\ref{sec:algorithms} presents algorithms. Section~\ref{sec:examples}
collects worked examples. Section~\ref{sec:experiments} reports experiments.
Sections~\ref{sec:discussion} and~\ref{sec:conclusion} discuss limitations and
conclude. Appendices collect proofs, notation, algorithmic details, and expanded
examples.


%% ══════════════════════════════════════════════════════════════════════════════
\section{Linear Algebra Foundations for Spectral Methods}
\label{sec:linalg}
%% ══════════════════════════════════════════════════════════════════════════════

This section establishes the linear-algebraic language used throughout the
paper. A reader comfortable with eigendecompositions and quadratic forms may
treat this section as a reference.

\subsection{Vectors as Signals}

Let $V = \{v_1,\dots,v_n\}$ be a finite set of $n$ entities. A \emph{signal}
on $V$ is a function $x\colon V \to \R$, represented as a column vector
$x \in \R^n$ with $x_i = x(v_i)$. We equip $\R^n$ with the standard inner
product $\langle x, y\rangle = x^\top y$ and induced norm $\|x\| =
\sqrt{x^\top x}$. Signals are the basic objects on which all operators act.

\subsection{Matrices as Operators}

A matrix $A \in \R^{n\times n}$ acts on signals by $x \mapsto Ax$. The $(i,j)$
entry of $A$ records how much the $j$-th coordinate of the input influences the
$i$-th coordinate of the output. Composition of operators corresponds to matrix
multiplication.

\subsection{Symmetry and Positive Semidefiniteness}

A matrix $A$ is \emph{symmetric} if $A = A^\top$. Symmetric matrices satisfy
$x^\top Ay = y^\top Ax$ for all $x,y$, which is the self-adjoint property. All
operators constructed in this paper will be symmetric.

\begin{definition}[Positive Semidefiniteness]
A symmetric matrix $A \in \R^{n\times n}$ is \emph{positive semidefinite}
(PSD), written $A \succeq 0$, if
\[
  x^\top A x \;\geq\; 0 \qquad \text{for all } x \in \R^n.
\]
It is \emph{positive definite} if $x^\top Ax > 0$ for all $x \neq 0$.
\end{definition}

Every PSD matrix can be written as $A = M^\top M$ for some (not necessarily
square) matrix $M$. This factored form is the key reason all Hodge Laplacians
are PSD: they are defined precisely as sums of such products.

\subsection{Quadratic Forms and Energy}

For a symmetric matrix $A$, the function $\mathcal{E}_A(x) = x^\top Ax$ is the
\emph{quadratic form} of $A$. It assigns a scalar ``energy'' to each signal
$x$. When $A \succeq 0$, this energy is nonnegative, and signals with zero
energy lie in the \emph{nullspace} $\ker(A) = \{x : Ax = 0\}$. The nullspace
is a linear subspace encoding the operator's equilibrium configurations.

\subsection{Energy as Smoothness and Structure}

A quadratic form of the type $x^\top A x$ should not be viewed as an abstract algebraic object, but as a measure of structural compatibility of a signal $x$ with respect to an operator $A$.

In the context of graph-based operators, this quantity admits a concrete interpretation. Let $L$ denote a graph Laplacian. Then one can expand:
\[
x^\top L x = \sum_{(u,v) \in E} w_{uv}(x_u - x_v)^2.
\]

This expression reveals that the quadratic form penalizes variation across edges: it is small when adjacent nodes take similar values, and large when they differ significantly. Thus, $x^\top L x$ measures the \emph{roughness} of the signal $x$ over the graph.

\paragraph{Convention on the energy sum.}
If one sums over unordered edges $\{u,v\}\in E$ once, then
$x^\top Lx = \sum_{\{u,v\}\in E} w_{uv}(x_u-x_v)^2$.
If one sums over ordered pairs via the adjacency matrix, the symmetric form
$x^\top Lx = \tfrac{1}{2}\sum_{u,v}A_{uv}(x_u-x_v)^2$ from Proposition~\ref{prop:graph_energy} is the same quantity. We use whichever is notationally convenient.

Eigenvectors of $L$ arise as minimizers of this energy under orthogonality constraints. The smallest eigenvalues correspond to the smoothest signals. In particular, the second eigenvector (the Fiedler vector) identifies the weakest point of connectivity in the graph, separating it into coherent regions.

This variational perspective will generalize directly to the higher-order and multi-layer operators constructed later in the paper.

\subsection{Eigenvalues, Eigenvectors, and Orthogonality}

\begin{definition}[Eigendecomposition]
A nonzero vector $u \in \R^n$ and scalar $\lambda \in \R$ satisfying
\[
  A u = \lambda u
\]
are an \emph{eigenvector} and \emph{eigenvalue} of $A$, respectively.
\end{definition}

The \emph{spectral theorem} guarantees that any symmetric $A$ has a complete
orthonormal eigenvector basis $u_1,\dots,u_n$ with real eigenvalues
$\lambda_1 \leq \lambda_2 \leq \cdots \leq \lambda_n$. Writing
$A = U\Lambda U^\top$ (with $U = [u_1 \cdots u_n]$ and $\Lambda =
\diag(\lambda_1,\dots,\lambda_n)$) makes the energy transparent:
$x^\top A x = \sum_i \lambda_i (u_i^\top x)^2$, a nonneg combination of
projected energies along each eigenmode.

\subsection{Rayleigh Quotients and Variational Thinking}

\begin{definition}[Rayleigh Quotient]
For symmetric $A$ and nonzero $x \in \R^n$, the \emph{Rayleigh quotient} is
\[
  R_A(x) = \frac{x^\top A x}{x^\top x}.
\]
\end{definition}

The variational principle states $\lambda_1 = \min_{x \neq 0} R_A(x)$, achieved
at $x = u_1$. More generally, $\lambda_j = \min\{R_A(x) : x \perp
u_1,\dots,u_{j-1},\, x\neq 0\}$. Eigenvectors are therefore the directions of
extremal energy, and eigenvalues are the corresponding energy levels.

\begin{example}[Two-by-two spectral example]
\label{ex:2x2}
Let
\[
  A = \begin{bmatrix} 2 & -1 \\ -1 & 2 \end{bmatrix}.
\]
The characteristic polynomial is $\lambda^2 - 4\lambda + 3 = 0$, giving
$\lambda_1 = 1$ and $\lambda_2 = 3$ with eigenvectors
$u_1 = \tfrac{1}{\sqrt{2}}[1,1]^\top$ and
$u_2 = \tfrac{1}{\sqrt{2}}[1,-1]^\top$. The in-phase direction $(u_1)$ incurs
energy $1$ while the out-of-phase direction $(u_2)$ incurs energy $3$: the
operator rewards agreement between coordinates and penalizes disagreement. This
$2\times 2$ template is the algebraic skeleton of every operator in this paper.
\end{example}


%% ══════════════════════════════════════════════════════════════════════════════
\section{Graphs and the Classical Graph Laplacian}
\label{sec:graph}
%% ══════════════════════════════════════════════════════════════════════════════

This section builds the first and most fundamental spectral operator on
relational data: the graph Laplacian. All operators introduced later reduce to
this one in the appropriate special case.

\subsection{Weighted Graphs and Adjacency Structure}

\begin{definition}[Weighted Graph]
A \emph{weighted undirected graph} is a triple $G = (V, E, w)$ where
$V = \{v_1,\dots,v_n\}$ is the vertex set, $E \subseteq \binom{V}{2}$ is the
edge set, and $w\colon E \to \R_{>0}$ assigns positive weights. The
\emph{adjacency matrix} $A \in \R^{n\times n}$ has $A_{uv} = w(u,v)$ if
$\{u,v\}\in E$ and $A_{uv}=0$ otherwise.
\end{definition}

\subsection{Degree Matrices and the Operator $L = D-A$}

\begin{definition}[Degree Matrix and Graph Laplacian]
The \emph{degree matrix} $D = \diag(d_1,\dots,d_n)$ has
\[
  d_u = \sum_v A_{uv}.
\]
The \emph{graph Laplacian} is $L = D - A$.
\end{definition}

Every row of $L$ sums to zero ($L\mathbf{1}=0$), reflecting that a constant
signal has zero energy on any graph. The matrix $L$ is symmetric because $A$ is
symmetric.

\subsection{Interpretation of the Graph Laplacian}

The graph Laplacian is not merely an algebraic construction, but a discrete differential operator encoding how signals vary across a network.

The identity
\[
x^\top L x = \sum_{(u,v) \in E} w_{uv}(x_u - x_v)^2
\]
shows that the Laplacian measures disagreement between neighboring nodes. A signal that is constant across connected components has zero energy and lies in the kernel of $L$.

The eigenvalues of $L$ quantify connectivity. In particular:
\begin{itemize}
    \item The multiplicity of the eigenvalue $0$ equals the number of connected components.
    \item The second-smallest eigenvalue (the Fiedler value) measures how well-connected the graph is.
    \item The corresponding eigenvector identifies a minimal-energy partition of the graph.
\end{itemize}

This interpretation will serve as the foundation for understanding higher-order Laplacians and their spectral behavior.

\subsection{Smoothness Energy on Graphs}

\begin{proposition}[Graph Laplacian Energy Identity]
\label{prop:graph_energy}
For any signal $x \in \R^n$,
\[
  x^\top L x = \frac{1}{2} \sum_{u,v} A_{uv}(x_u - x_v)^2.
\]
\end{proposition}

\begin{proof}
We expand $x^\top(D-A)x = x^\top Dx - x^\top Ax = \sum_u d_u x_u^2 -
\sum_{u,v}A_{uv}x_u x_v$. Using $d_u = \sum_v A_{uv}$ in the first term:
\[
  x^\top Lx = \sum_{u,v}A_{uv}x_u^2 - \sum_{u,v}A_{uv}x_u x_v
  = \tfrac{1}{2}\sum_{u,v}A_{uv}(x_u^2 - 2x_u x_v + x_v^2)
  = \tfrac{1}{2}\sum_{u,v}A_{uv}(x_u-x_v)^2. \qedhere
\]
\end{proof}

The energy $x^\top Lx$ is a weighted sum of squared differences across all
edges. A signal with low energy varies slowly across highly-weighted edges; this
is the formal meaning of \emph{graph smoothness}. Since every term is
nonnegative, $L\succeq 0$, and the eigenvalues satisfy
\[
  0 = \lambda_1 \leq \lambda_2 \leq \cdots \leq \lambda_n.
\]

\subsection{Kernel, Connectedness, and Constant Signals}

\begin{proposition}[Kernel of the Graph Laplacian]
\label{prop:graph_kernel}
$x \in \ker(L)$ if and only if $x_u = x_v$ for every edge $\{u,v\}\in E$.
Consequently: (1) constant signals always lie in $\ker(L)$; (2) for a
connected graph, $\ker(L) = \spn(\mathbf{1})$ and $\lambda_1=0$ has multiplicity
one; (3) for a graph with $c$ connected components, $\dim\ker(L) = c$.
\end{proposition}

\begin{proof}
By Proposition~\ref{prop:graph_energy}, $x^\top Lx = 0$ iff $A_{uv}(x_u-x_v)^2
= 0$ for all $u,v$, i.e., $x_u = x_v$ on every edge. Connectedness forces
every such $x$ to be constant. The component count follows by applying the
same argument within each component.
\end{proof}

Constant signals represent perfect global consensus: every vertex holds the
same value. The nullity of $L$ counts how many independent consensus modes
exist---one per connected component.

\subsection{The Second Eigenvalue and the Fiedler Vector}

\begin{definition}[Fiedler Value and Vector]
For a connected graph, $\lambda_2 > 0$ is the \emph{Fiedler value} (algebraic
connectivity). The corresponding unit eigenvector $u_2 \perp \mathbf{1}$ is the
\emph{Fiedler vector}.
\end{definition}

The Fiedler vector partitions vertices by sign: vertices with $u_2(v)>0$ form
one cluster, those with $u_2(v)<0$ form another~\cite{chung1997}. A large
$\lambda_2$ indicates high algebraic connectivity (the graph is hard to cut);
a small $\lambda_2$ flags a near-bottleneck. The Fiedler theory generalizes
directly to the Supra-Hodge setting in Section~\ref{sec:spectral}.

\subsection{Worked Graph Example}

\begin{example}[Four-node weighted graph]
\label{ex:4node}
Let $G$ have $V=\{1,2,3,4\}$ and edges $\{1,2\},\{2,3\},\{3,4\},\{4,1\},\{1,3\}$,
all with unit weight. Then
\begingroup
\setlength{\arraycolsep}{3pt}
\small
\[
  \begin{aligned}
    A &= \begin{bmatrix}
      0 & 1 & 1 & 1\\
      1 & 0 & 1 & 0\\
      1 & 1 & 0 & 1\\
      1 & 0 & 1 & 0
    \end{bmatrix},\\[4pt]
    D &= \diag(3,2,3,2),\\[4pt]
    L &= \begin{bmatrix}
      3 & -1 & -1 & -1\\
      -1 & 2 & -1 & 0\\
      -1 & -1 & 3 & -1\\
      -1 & 0 & -1 & 2
    \end{bmatrix}.
  \end{aligned}
\]
\endgroup
The constant signal $x=\mathbf{1}$ satisfies $Lx=0$. The signal
$y=[1,-1,1,-1]^\top$ has energy $y^\top Ly = 8$, reflecting large differences
across all edges. The Fiedler vector separates $\{1,4\}$ from $\{2,3\}$,
marking the minimal bottleneck. We reuse this example in
Section~\ref{sec:examples}.
\end{example}


%% ══════════════════════════════════════════════════════════════════════════════
\section{Multilayer Graphs and Supra-Laplacians}
\label{sec:multilayer}
%% ══════════════════════════════════════════════════════════════════════════════

This section extends the graph Laplacian to systems with multiple relation
types over the same entity set. The resulting supra-Laplacian is the direct
predecessor of the Supra-Hodge Laplacian, and its energy identity is the
template for the main variational result of the paper.

\subsection{Multiple Relation Types on One Entity Set}

Let $V$ be a fixed entity set with $|V|=n$. We consider $k$ weighted graphs
$\mathcal{G} = \{G^{(1)},\dots,G^{(k)}\}$, each on vertex set $V$:
\[
  \mathcal{G} = \{G^{(1)}, G^{(2)}, \dots, G^{(k)}\}.
\]
Each $G^{(i)} = (V, E^{(i)}, w^{(i)})$ encodes a distinct relation type. In a
reasoning system, these might be semantic similarity, evidential support,
causal dependency, and temporal precedence.

\subsection{Layerwise Graph Laplacians}

Each layer $G^{(i)}$ has its own adjacency matrix $A^{(i)}$, degree matrix
$D^{(i)}$, and graph Laplacian $L^{(i)} = D^{(i)} - A^{(i)}$. By
Proposition~\ref{prop:graph_energy}, $L^{(i)} \succeq 0$ and $\ker(L^{(i)})$
captures the consensus structure within layer $i$.

\subsection{Stacked Signals and Block Operators}

To analyze all layers jointly, we stack layer signals into a single vector
$v = (v^{(1)},\dots,v^{(k)}) \in \R^{kn}$. Inter-layer coupling is encoded by
weights $\omega_{ij} \geq 0$: a positive $\omega_{ij}$ attracts the signals of
layers $i$ and $j$ toward agreement.

\subsection{The Supra-Laplacian and Its Energy}

\begin{definition}[Supra-Laplacian]
The \emph{supra-Laplacian} is the $kn\times kn$ block-symmetric matrix
\begin{equation}
  \label{eq:supra}
  \Lsup =
  \begin{bmatrix}
    L^{(1)} + \textstyle\sum_{j\neq 1}\omega_{1j}I & -\omega_{12}I & \cdots & -\omega_{1k}I\\
    -\omega_{21}I & L^{(2)} + \textstyle\sum_{j\neq 2}\omega_{2j}I & \cdots & -\omega_{2k}I\\
    \vdots & \vdots & \ddots & \vdots\\
    -\omega_{k1}I & -\omega_{k2}I & \cdots & L^{(k)} + \textstyle\sum_{j\neq k}\omega_{kj}I
  \end{bmatrix}.
\end{equation}
\end{definition}

\begin{proposition}[Supra-Laplacian Energy Decomposition]
\label{prop:supra_energy}
For any stacked signal $v=(v^{(1)},\dots,v^{(k)})$,
\begingroup
\small
\[
  \begin{aligned}
    v^\top \Lsup v
    &= \sum_{i=1}^k v^{(i)\top} L^{(i)} v^{(i)} \\
    &\quad + \sum_{i<j}\omega_{ij}\|v^{(i)} - v^{(j)}\|_2^2.
  \end{aligned}
\]
\endgroup
\end{proposition}

\begin{proof}
Each diagonal block $L^{(i)}+\sum_{j\neq i}\omega_{ij}I$ contributes
$v^{(i)\top}L^{(i)}v^{(i)} + \sum_{j\neq i}\omega_{ij}\|v^{(i)}\|^2$ to
$v^\top\Lsup v$. Each off-diagonal block $-\omega_{ij}I$ contributes
$-2\omega_{ij}v^{(i)\top}v^{(j)}$. Collecting all pairs $(i,j)$ with $i<j$:
\[
  \sum_{i<j}\omega_{ij}\bigl(\|v^{(i)}\|^2 - 2v^{(i)\top}v^{(j)}
  + \|v^{(j)}\|^2\bigr)
  = \sum_{i<j}\omega_{ij}\|v^{(i)}-v^{(j)}\|^2. \qedhere
\]
\end{proof}

The energy decomposes into within-layer graph smoothness and cross-layer
disagreement. This structure is the exact prototype for the Supra-Hodge energy
expansion proved in Section~\ref{sec:properties}.

\subsection{Consensus and Disagreement Across Layers}

A stacked signal $v$ achieves zero supra-Laplacian energy iff $L^{(i)}v^{(i)}=0$
for all $i$ (within-layer graph consensus) and $v^{(i)}=v^{(j)}$ whenever
$\omega_{ij}>0$ (cross-layer consensus). The supra-Laplacian therefore enforces
two simultaneous smoothness conditions.

\subsection{Worked Two-Layer Example}

\begin{example}[Two-layer supra-Laplacian]
\label{ex:supra}
Let $n=3$, $k=2$, both layers equal to the path graph $1$--$2$--$3$ with
$L^{(1)}=L^{(2)}=L_{\mathrm{path}}$, and $\omega_{12}=1$. Then
\[
  \Lsup = \begin{bmatrix} L_{\mathrm{path}}+I & -I \\ -I & L_{\mathrm{path}}+I \end{bmatrix}.
\]
The symmetric signal $v=(\mathbf{1},\mathbf{1})$ has zero energy (constant in
each layer, equal across layers). The antisymmetric signal
$v=(\mathbf{1},-\mathbf{1})$ incurs cross-layer energy
$\omega_{12}\|2\mathbf{1}\|^2=12$.
\end{example}

\subsection{Why Pairwise Multilayer Models Still Fall Short}

Even with $k$ coupled layers, the supra-Laplacian encodes only pairwise edges
within each layer. It cannot distinguish a filled triangle (a 2-simplex) from a
hollow cycle of three edges. If three entities jointly support a conclusion in a
way that no pair does independently, the supra-Laplacian assigns the same energy
to both configurations. Correcting this requires Hodge theory.


%% ══════════════════════════════════════════════════════════════════════════════
\section{Simplicial Complexes and Higher-Order Relational Structure}
\label{sec:simplicial}
%% ══════════════════════════════════════════════════════════════════════════════

This section introduces simplicial complexes: the combinatorial structures that
generalize graphs to encode interactions involving arbitrarily many entities
simultaneously.

\subsection{From Graphs to Simplicial Complexes: A Concrete Example}

A graph encodes pairwise relationships through edges. However, many systems exhibit interactions that involve more than two entities simultaneously.

Consider three vertices $u, v, w$. A graph may contain all three edges $(u,v)$, $(v,w)$, and $(u,w)$, forming a triangle. However, this does not necessarily imply that the three entities form a coherent higher-order unit.

A simplicial complex distinguishes between these cases:
\begin{itemize}
    \item The presence of all three edges represents pairwise compatibility.
    \item The inclusion of the 2-simplex $(u,v,w)$ asserts a genuine three-way interaction.
\end{itemize}

This distinction is critical: higher-order structure cannot be inferred from pairwise relations alone. Simplicial complexes provide the appropriate mathematical language for encoding such multi-way interactions.

\subsection{From Pairwise Edges to Simplices}

\begin{definition}[Simplex and Simplicial Complex]
A \emph{$p$-simplex} is a set $\sigma = \{v_0,v_1,\dots,v_p\} \subseteq V$ of
$p+1$ elements. A \emph{simplicial complex} $K$ on $V$ is a collection of
simplices satisfying the \emph{closure property}:
\[
  \sigma \in K,\quad \tau \subseteq \sigma \implies \tau \in K.
\]
\end{definition}

The dimension of a simplex is $\dim(\sigma) = |\sigma|-1$. Thus 0-simplices
are vertices, 1-simplices are (undirected) edges, 2-simplices are filled
triangles, and 3-simplices are filled tetrahedra.

\subsection{Faces, Closure, and Dimension}

Each subset of a simplex is a \emph{face}. The closure property says: if a
higher-order interaction is present, all its sub-interactions are present too.
This is the natural consistency requirement for relational data. We write
$K_p$ for the set of all $p$-simplices of $K$.

\subsection{Hollow Triangles Versus Filled Triangles}

\begin{example}[Hollow vs.\ filled triangle]
\label{ex:hollow}
Let $V=\{1,2,3\}$. The \emph{hollow complex}
$K_{\mathrm{hollow}} = \{\{1\},\{2\},\{3\},\{1,2\},\{1,3\},\{2,3\}\}$ contains
all three edges but not the triangle as a 2-simplex. The \emph{filled complex}
$K_{\mathrm{filled}} = K_{\mathrm{hollow}} \cup \{\{1,2,3\}\}$ adds the triangle.
These two complexes have identical pairwise adjacency structure. Yet they have
different topology: $K_{\mathrm{hollow}}$ contains a non-trivial 1-cycle (a
``hole''), while $K_{\mathrm{filled}}$ does not. Graph Laplacian methods cannot
distinguish them; Hodge methods detect the difference precisely via harmonic
structure.
\end{example}

\subsection{Why Higher-Order Structure Changes the Model Class}

Adding the 2-simplex $\{1,2,3\}$ changes the topology of the complex and
changes the spectrum of the associated Hodge Laplacian. In the hollow complex,
the edge signal $f=[1,1,1]^\top$ (uniform flow around the cycle) is harmonic
at $p=1$ (zero $H_1$-energy). In the filled complex, this signal has nonzero
upper Laplacian energy: the triangle ``detects'' that the flow winds around it,
and the harmonic cycle is destroyed. No pairwise operator can see this
difference.

\subsection{Examples from Reasoning and Relational Systems}

In an evidential layer of a reasoning trace, three propositions
$\{P_1, P_2, P_3\}$ may jointly support a conclusion even though no pair does
alone. This is a 2-simplex in the evidential complex. Its presence eliminates
the harmonic 1-cycle formed by the three pairwise support edges, flagging the
triplet as genuinely consistent. Pairwise models would treat a three-way
evidence coalition identically to three isolated pairwise connections.


%% ══════════════════════════════════════════════════════════════════════════════
\section{Chain Spaces, Orientation, and Boundary Operators}
\label{sec:chains}
%% ══════════════════════════════════════════════════════════════════════════════

This section develops the algebraic machinery for computing with simplicial
complexes: chain spaces that coordinatize signals on simplices of each
dimension, and boundary operators that map between consecutive dimensions.

\subsection{$p$-Simplices and $p$-Chains}

\begin{definition}[Chain Space]
The \emph{$p$-chain space} of $K$ is
\[
  C_p(K) = \R^{|K_p|},
\]
with one coordinate per $p$-simplex of $K$. A \emph{$p$-chain} is an element
of $C_p(K)$.
\end{definition}

Signals on vertices live in $C_0(K)=\R^n$; signals on edges live in $C_1(K)$;
signals on triangles live in $C_2(K)$. The chain space generalizes the notion
of a vertex signal to signals of every order.

\subsection{Basis Elements and Coordinate Representations}

We choose an orientation for each $p$-simplex. An oriented simplex is written
$[v_0,v_1,\dots,v_p]$, and reversing any adjacent pair negates the element:
$[v_0,\dots,v_i,v_{i+1},\dots,v_p] = -[v_0,\dots,v_{i+1},v_i,\dots,v_p]$.
The set of consistently oriented $p$-simplices forms a basis for $C_p(K)$.

\subsection{Orientation and Sign}

Orientation is a bookkeeping device that makes boundary operators well-defined.
For an edge $[u,v]$, it assigns a direction (arrow from $u$ to $v$). For a
triangle $[u,v,w]$, it assigns a cyclic ordering. Changing orientation changes
the sign of the basis element, ensuring that boundary operators are
orientation-consistent.

\subsection{Boundary Operators}

\begin{definition}[Boundary Operator]
The \emph{boundary operator} $\partial_p\colon C_p(K)\to C_{p-1}(K)$ is
defined on basis elements by
\[
  \partial_p[v_0,v_1,\dots,v_p]
  = \sum_{i=0}^{p}(-1)^i[v_0,\dots,\hat{v}_i,\dots,v_p],
\]
where $\hat{v}_i$ denotes omission of $v_i$, and extended by linearity.
\end{definition}

\subsection{Example: Boundary of a Triangle}

For a 2-simplex $[v_0,v_1,v_2]$, one has
$\partial_2[v_0,v_1,v_2]=[v_1,v_2]-[v_0,v_2]+[v_0,v_1]$, the oriented edge cycle
around the triangle; applying $\partial_1$ again gives $0$, i.e.\ $\partial_1\partial_2=0$
(matrix form: $B_1B_2=0$). Example~\ref{ex:boundary} computes the incidence matrices
explicitly.

The boundary operator maps each oriented $p$-simplex to the signed sum of its
$(p-1)$-dimensional faces, with signs determined by the orientation. Intuitively,
the boundary of an oriented edge $[u,v]$ is the signed difference $v-u$; the
boundary of an oriented triangle $[u,v,w]$ is the signed cycle
$[v,w]-[u,w]+[u,v]$.

\subsection{Incidence Matrices}

We represent $\partial_p$ as the matrix $B_p \in \R^{|K_{p-1}|\times|K_p|}$:
the column for oriented $p$-simplex $\sigma$ has $+1$ in row $\tau$ if $\tau$
appears in $\partial_p\sigma$ with positive orientation, $-1$ if with negative
orientation, and $0$ otherwise.

\subsection{Why Boundary of Boundary Is Zero}

\begin{proposition}[Boundary of Boundary Identity]
\label{prop:b2}
For any $p$,
\[
  \partial_p\partial_{p+1} = 0, \qquad B_pB_{p+1} = 0.
\]
\end{proposition}

\begin{proof}
For an oriented $(p+1)$-simplex $[v_0,\dots,v_{p+1}]$, compute
$\partial_p\partial_{p+1}[v_0,\dots,v_{p+1}]$:
\[
  \partial_p\partial_{p+1}[v_0,\dots,v_{p+1}]
  = \sum_{i=0}^{p+1}(-1)^i\partial_p[v_0,\dots,\hat{v}_i,\dots,v_{p+1}].
\]
Each $p$-face $[v_0,\dots,\hat{v}_i,\dots,\hat{v}_j,\dots,v_{p+1}]$ (with both
$v_i$ and $v_j$ omitted) appears twice: once as the $j$-th term of
$\partial_p[v_0,\dots,\hat{v}_i,\dots]$ with sign $(-1)^i(-1)^{j-1}$, and
once as the $i$-th term of $\partial_p[v_0,\dots,\hat{v}_j,\dots]$ with sign
$(-1)^j(-1)^i$. These two contributions are equal in magnitude and opposite in
sign, so they cancel. Hence $\partial_p\partial_{p+1}=0$. The matrix identity
$B_pB_{p+1}=0$ is the coordinate form.
\end{proof}

This identity is the algebraic cornerstone of homological algebra. It implies
$\im(B_{p+1})\subseteq\ker(B_p)$, enabling the Hodge decomposition.

\subsection{Worked Simplex-Boundary Example}

\begin{example}[Oriented triangle boundary]
\label{ex:boundary}
Let $K$ have vertices $\{1,2,3\}$ and oriented edges
$[1,2],[1,3],[2,3]$. Let the oriented triangle be $[1,2,3]$. Then
\[
  \partial_2[1,2,3] = [2,3] - [1,3] + [1,2].
\]
With edges listed in order $[1,2],[1,3],[2,3]$ as the basis of $C_1$:
\[
  B_2 = \begin{bmatrix} 1 \\ -1 \\ 1 \end{bmatrix},\quad
  B_1 = \begin{bmatrix} -1 & -1 & 0 \\ 1 & 0 & -1 \\ 0 & 1 & 1 \end{bmatrix}.
\]
Direct computation confirms $B_1 B_2 = [1-1, -1+1, 1-1]^\top = 0$,
illustrating Proposition~\ref{prop:b2}.
\end{example}


%% ══════════════════════════════════════════════════════════════════════════════
\section{Hodge Laplacians on Simplicial Complexes}
\label{sec:hodge}
%% ══════════════════════════════════════════════════════════════════════════════

We now define the combinatorial Hodge Laplacian, the higher-order spectral
operator on a single complex that the Supra-Hodge Laplacian will generalize.

\subsection{Definition of the Combinatorial Hodge Laplacian}

\begin{definition}[Hodge Laplacian]
The \emph{$p$-th combinatorial Hodge Laplacian} of $K$ is
\[
  H_p = B_p^\top B_p + B_{p+1}B_{p+1}^\top,
\]
where we set $B_0\equiv 0$ by convention. The term $B_p^\top B_p$ is the
\emph{lower Laplacian} and $B_{p+1}B_{p+1}^\top$ is the \emph{upper Laplacian}.
\end{definition}

\subsection{Interpretation of Hodge Decomposition}

By the standard combinatorial Hodge decomposition for finite simplicial
complexes~\cite{horak2013,schaub2020,lim2020}, chain space splits into three
orthogonal components:
\[
C_p = \mathrm{im}(B_p^\top) \oplus \ker(H_p) \oplus \mathrm{im}(B_{p+1}).
\]

These components admit intuitive interpretations:
\begin{itemize}
    \item \textbf{Gradient component:} signals derived from lower-order structure.
    \item \textbf{Curl component:} signals corresponding to local cycles or inconsistencies.
    \item \textbf{Harmonic component:} globally consistent structures that are neither gradients nor curls.
\end{itemize}

In particular, for $p = 1$, the curl term captures inconsistencies on triangles and higher simplices. This enables the detection of cyclic contradictions that are invisible to graph-based methods.

The lower Laplacian measures how much a $p$-chain signal diverges at
$(p-1)$-boundaries; the upper Laplacian measures how much it has nonzero curl
around $(p+1)$-simplices.

\begin{proposition}[Hodge Laplacians Are Symmetric PSD]
\label{prop:hodge_psd}
$H_p$ is symmetric and positive semidefinite for every $p$.
\end{proposition}

\begin{proof}
Both $B_p^\top B_p$ and $B_{p+1}B_{p+1}^\top$ are of the form $M^\top M$
(with $M=B_p$ and $M=B_{p+1}^\top$ respectively), hence symmetric PSD. A sum
of symmetric PSD matrices is symmetric PSD.
\end{proof}

\subsection{Recovery of the Graph Laplacian at $p=0$}

\begin{proposition}[Hodge-0 Equals Graph Laplacian]
\label{prop:h0}
$H_0 = B_1 B_1^\top$ equals the graph Laplacian $L$ of the 1-skeleton of $K$.
\end{proposition}

\begin{proof}
For oriented edge $[u,v]$ with $u<v$ (so that $B_1$ has $-1$ in row $u$ and
$+1$ in row $v$), the $(a,b)$-entry of $B_1 B_1^\top$ is
$\sum_e (B_1)_{ae}(B_1)_{be}$. For $a=b$: $(B_1 B_1^\top)_{aa} = \sum_e
(B_1)_{ae}^2 = |\{e : a \in e\}| = d_a$. For $a\neq b$: $(B_1B_1^\top)_{ab}
= -A_{ab}$ (nonzero only when $\{a,b\}\in E$). Hence $B_1B_1^\top = D-A = L$.
\end{proof}

This proposition confirms that Hodge theory is a genuine extension of graph
spectral theory: $H_0$ is the original graph Laplacian.

\subsection{The $p=1$ Case and Cycle-Level Inconsistency}

For $p=1$, a signal $f\in C_1(K)$ (a flow or relation strength on edges) has
\[
  f^\top H_1 f = \|B_1 f\|^2 + \|B_2^\top f\|^2.
\]
The term $\|B_1 f\|^2$ measures vertex-level divergence of $f$ (how much flow
is not conserved at each vertex). The term $\|B_2^\top f\|^2$ measures
triangle-level curl (how much $f$ winds around each filled triangle). Both
terms must vanish for $f$ to be harmonic.

\subsection{Lower and Upper Adjacency}

Two $p$-simplices are \emph{lower-adjacent} if they share a $(p-1)$-face; they
are \emph{upper-adjacent} if they are co-faces of a common $(p+1)$-simplex.
Lower adjacency governs the lower Laplacian $B_p^\top B_p$; upper adjacency
governs the upper Laplacian $B_{p+1}B_{p+1}^\top$. Graph Laplacians correspond
to the $p=0$ case where only lower adjacency (vertex-to-edge incidence) exists.

\subsection{Gradient, Curl, and Harmonic Components}

The Hodge decomposition theorem states that for any $f\in C_p(K)$,
\[
  f = f_{\mathrm{grad}} + f_{\mathrm{harm}} + f_{\mathrm{curl}},
\]
where $f_{\mathrm{grad}}\in\im(B_p^\top)$, $f_{\mathrm{curl}}\in\im(B_{p+1})$,
and $f_{\mathrm{harm}}\in\ker(H_p)$. These subspaces are mutually orthogonal.
The gradient component is the conservative part (derivable from a vertex
potential); the curl component is the exact part (boundary of a higher-order
chain); the harmonic component is the topologically persistent part (belonging
to genuine higher-order cycles).

\paragraph{Remark on the decomposition.}
For finite simplicial complexes equipped with the standard inner product on chain spaces, the combinatorial Hodge decomposition is an orthogonal decomposition. Thus the gradient, harmonic, and curl components are not merely formal labels, but mutually orthogonal subspaces with distinct spectral roles.

\subsection{Hodge Decomposition Intuition}

On edge signals ($p=1$), the gradient component represents potential-driven flow
(water flowing downhill from a potential), the curl component represents
circulation induced by filled triangles (water flowing around a filled basin),
and the harmonic component represents persistent circulation that is not caused
by any local source or triangle fill (water in a donut). The Hodge Laplacian
measures and separates all three simultaneously.

\subsection{Worked Hodge Example}

\begin{example}[Filled-triangle Hodge computation]
\label{ex:hodge}
Let $K=K_{\mathrm{filled}}$ from Example~\ref{ex:hollow}, with
$B_1$ and $B_2$ as in Example~\ref{ex:boundary}. Then
\[
  B_1^\top B_1 = \begin{bmatrix}2&1&-1\\1&2&1\\-1&1&2\end{bmatrix},\quad
  B_2 B_2^\top = \begin{bmatrix}1&-1&1\\-1&1&-1\\1&-1&1\end{bmatrix},
\]
giving
\[
  H_1 = B_1^\top B_1 + B_2B_2^\top
      = \begin{bmatrix}3&0&0\\0&3&0\\0&0&3\end{bmatrix}.
\]
Thus $H_1=3I_3$ for the filled triangle, and $\ker(H_1)=\{0\}$: there are no
harmonic 1-cycles. For the hollow complex ($B_2=0$), $H_1=B_1^\top B_1$ has
$\ker(H_1) = \spn\bigl([1,1,1]^\top/\sqrt{3}\bigr)$---the uniform circulation
around the empty triangle is a harmonic cycle.
\end{example}


%% ══════════════════════════════════════════════════════════════════════════════
\section{From Single Higher-Order Systems to Coupled Higher-Order Systems}
\label{sec:coupling}

A single simplicial complex captures one relational view. Coupling several such
views requires a common ambient $p$-chain space so that layerwise signals and
operators can be compared and coupled without conflating distinct per-view
topologies.

The purpose of this section is to construct the aligned ambient chain space and
embedding machinery needed to couple multiple higher-order views. Once that
infrastructure is in place, the Supra-Hodge operator itself will be defined
formally in Section~\ref{sec:defn}.

\subsection{Why Alignment Is Necessary}

Each layer in a multi-view system defines its own simplicial complex and corresponding chain space. Without a common coordinate system, these spaces cannot be meaningfully coupled.

A naive approach would attempt to compare signals across layers directly. However, if the underlying simplices differ, such comparisons are ill-defined: one layer may contain a simplex that does not exist in another.

To resolve this, we introduce an aligned ambient chain space. Each layer is embedded into this shared space by zero-padding missing simplices. This construction enables the definition of a block operator that couples corresponding structures across layers while preserving the internal topology of each layer.

This alignment step is the key technical device that allows higher-order operators to be extended to the multi-layer setting.

\subsection{Why a Single Simplicial Complex Is Not Enough}

A reasoning agent maintains semantic clusters, evidential support relations, and
temporal precedences over the same entities. Each view has its own higher-order
structure (semantic triplets, evidence triplets, temporal chains). Merging all
views into a single complex conflates semantically unrelated triangles and
destroys the per-view interpretation of harmonic cycles. We need a family of
complexes and a principled coupling.

\subsection{Families of Complexes Over a Shared Entity Set}

\begin{definition}[Family of Simplicial Complexes]
Let $V$ be a shared vertex set. A \emph{family of simplicial complexes} is
\[
  \calK = \{K^{(1)}, K^{(2)}, \dots, K^{(k)}\},
\]
where each $K^{(i)}$ is a simplicial complex on $V$.
\end{definition}

Each $K^{(i)}$ has its own $p$-simplices $K_p^{(i)}$ and chain space
$C_p^{(i)} = \R^{|K_p^{(i)}|}$.

\subsection{Alignment of Simplex Coordinates Across Views}

To couple the chain spaces of different layers, we need a shared coordinate
system.

\begin{definition}[Aligned Reference Set]
An \emph{aligned reference set} $\Sigma_p$ is a set of admissible oriented
$p$-simplices such that $K_p^{(i)}\subseteq\Sigma_p$ for all $i$. The ambient
dimension is $N_p = |\Sigma_p|$. The \emph{common ambient chain space} is
\[
  \Ctil_p \;\cong\; \R^{N_p}.
\]
\end{definition}

The simplest choice is $\Sigma_p = \bigcup_i K_p^{(i)}$ (union over all
layers), with $N_p = |\bigcup_i K_p^{(i)}|$.

\subsection{Explicit Embedding into the Ambient Chain Space}

To make the alignment construction fully explicit, we define for each layer $i$ an embedding map
\[
E_p^{(i)} : C_p\!\left(K^{(i)}\right) \longrightarrow \Ctil_p \cong \R^{N_p},
\]
where $N_p = |\Sigma_p|$. The map $E_p^{(i)}$ sends each basis simplex of $C_p(K^{(i)})$
to its corresponding coordinate in the ambient basis indexed by $\Sigma_p$, and assigns zero to all coordinates corresponding to simplices absent from layer $i$.

Using this embedding, any layerwise $p$-chain signal $x^{(i)} \in C_p(K^{(i)})$ is represented in the common ambient space by $E_p^{(i)}x^{(i)} \in \Ctil_p$. After padding incidence matrices to the reference set $\Sigma_p$ (equivalently, conjugating the native Hodge matrices by $E_p^{(i)}$), every layerwise Hodge Laplacian becomes an $N_p\times N_p$ operator on $\Ctil_p$. We denote this embedded matrix by $H_p^{(i)}$ throughout the block definition \eqref{eq:Lp} and the remainder of the paper.

This makes the coupling mathematically precise: the Supra-Hodge operator does not compare signals in unrelated spaces, but in a shared ambient chain space obtained by explicit embeddings.

\paragraph{Notational convention.}
From this point onward, we write $H_p^{(i)}$ for the ambient-space operator acting on $\Ctil_p$ obtained from the embedding construction above. Thus all layerwise Hodge operators appearing in the definition of $L_p$, in the main theorem, and in subsequent reductions are understood as operators on the common aligned chain space.

\subsection{Zero-Padding and a Common Ambient Chain Space}

Each layer $i$ has its own chain space $C_p^{(i)}$ of dimension
$|K_p^{(i)}|\leq N_p$. We embed it into $\Ctil_p$ by \emph{zero-padding}: the
boundary matrix $B_p^{(i)}$ is extended to $\R^{N_{p-1}\times N_p}$ by
inserting zero columns for simplices in $\Sigma_p\setminus K_p^{(i)}$. The
Hodge Laplacian $H_p^{(i)} = B_p^{(i)\top}B_p^{(i)} + B_{p+1}^{(i)}
B_{p+1}^{(i)\top}$ then acts on the ambient space $\Ctil_p \cong \R^{N_p}$.

\begin{example}[Alignment and zero-padding]
\label{ex:align}
Let $k=2$, $V=\{1,2,3\}$, with $K_1^{(1)}=\{[1,2],[1,3]\}$ and
$K_1^{(2)}=\{[1,2],[2,3]\}$. Choose $\Sigma_1=\{[1,2],[1,3],[2,3]\}$,
so $N_1=3$. The aligned incidence matrix $B_1^{(1)}$ has a zero column for
$[2,3]$, and $B_1^{(2)}$ has a zero column for $[1,3]$. The Hodge Laplacians
$H_1^{(1)}$ and $H_1^{(2)}$ both live in $\R^{3\times 3}$ and can be assembled
into a $2\times 3 \cdot 2\times 3 = 6\times 6$ block operator.
\end{example}

\subsection{When Identity Coupling Is Natural}

If all layers share the same $p$-simplices ($K_p^{(i)}=K_p^{(j)}$ for all
$i,j$), then corresponding simplices are semantically paired and the natural
coupling operator between layers $i$ and $j$ is $C^{(ij)}=\omega_{ij}I_{N_p}$.
This is the scalar coupling case studied throughout this paper.

\subsection{When Identity Coupling Fails}

If layers have different complex structures with no natural simplex
correspondence, identity coupling may conflate semantically unrelated simplices.
In this case, a learned or structured coupling matrix $C^{(ij)}\in\R^{N_p\times
N_p}$ replaces $\omega_{ij}I_{N_p}$. This extension is discussed in
Section~\ref{sec:discussion}; we focus on the scalar case here.


%% ══════════════════════════════════════════════════════════════════════════════
\section{Definition of the Supra-Hodge Laplacian}
\label{sec:defn}
%% ══════════════════════════════════════════════════════════════════════════════

Having fixed the alignment, embedding, and notational conventions of
Section~\ref{sec:coupling}, we now define the Supra-Hodge Laplacian $L_p$, the
central operator of this paper.

\subsection{Layerwise Hodge Operators}

Given the family $\calK$ and aligned incidence matrices $B_p^{(i)}\in
\R^{N_{p-1}\times N_p}$, the layerwise Hodge Laplacian is
\[
  H_p^{(i)} = B_p^{(i)\top}B_p^{(i)} + B_{p+1}^{(i)}B_{p+1}^{(i)\top}
  \;\in\; \R^{N_p\times N_p}.
\]
By Proposition~\ref{prop:hodge_psd}, $H_p^{(i)}\succeq 0$. The aligned
construction ensures all $H_p^{(i)}$ act on the same ambient space $\Ctil_p$.

\subsection{Stacked $p$-Chain Signal Spaces}

The total signal space for the $k$-layer system at order $p$ is
\[
  \calC_p = \underbrace{\Ctil_p \oplus \cdots \oplus \Ctil_p}_{k}
  \;\cong\; \R^{kN_p}.
\]
A stacked signal $v=(v^{(1)},\dots,v^{(k)})\in\calC_p$ collects one ambient
$p$-chain signal from each layer.

\subsection{Inter-Layer Couplings and the Coupling Graph}

\begin{definition}[Coupling Graph and Coupling Operator]
The \emph{coupling graph} $\Gom$ has vertex set $\{1,\dots,k\}$ and edge
$\{i,j\}$ iff $\omega_{ij}>0$. The \emph{coupling operator} between layers $i$
and $j$ is $C^{(ij)}=\omega_{ij}I_{N_p}$.
\end{definition}

The coupling graph encodes which layer pairs are linked. Its connectivity
structure will determine the dimension of $\ker(L_p)$ through
Theorem~\ref{thm:main}.

\subsection{Block-Matrix Definition of $L_p$}

\begin{definition}[Supra-Hodge Laplacian]
\label{def:main}
The \emph{order-$p$ Supra-Hodge Laplacian} is the $kN_p\times kN_p$
block-symmetric matrix
\begin{equation}
  \label{eq:Lp}
  L_p =
  \begin{bmatrix}
    H_p^{(1)} + \textstyle\sum_{j\neq 1}\omega_{1j}I & -\omega_{12}I & \cdots & -\omega_{1k}I\\
    -\omega_{21}I & H_p^{(2)} + \textstyle\sum_{j\neq 2}\omega_{2j}I & \cdots & -\omega_{2k}I\\
    \vdots & \vdots & \ddots & \vdots\\
    -\omega_{k1}I & -\omega_{k2}I & \cdots & H_p^{(k)} + \textstyle\sum_{j\neq k}\omega_{kj}I
  \end{bmatrix}.
\end{equation}
\end{definition}

Each diagonal block $H_p^{(i)}+\sum_{j\neq i}\omega_{ij}I$ is the layerwise
Hodge operator augmented by the total coupling weight from layer $i$ to all
other layers. Each off-diagonal block $-\omega_{ij}I$ provides the cross-layer
attractive force. The structural analogy with the supra-Laplacian
equation~\eqref{eq:supra} is exact, with $H_p^{(i)}$ replacing $L^{(i)}$ and
$N_p$ replacing $n$.

\subsection{Equivalent Operator Form}

Let
\[
\mathcal{H}_p := \diag\!\bigl(H_p^{(1)},\dots,H_p^{(k)}\bigr)
\]
denote the block-diagonal operator collecting the layerwise Hodge Laplacians, and let $L_{\mathrm{inter}}$ denote the graph Laplacian of the coupling graph $\Gom$. Then the Supra-Hodge Laplacian may be written compactly as
\[
L_p = \mathcal{H}_p + L_{\mathrm{inter}} \otimes I_{N_p}.
\]

This form makes the structure of the operator especially transparent: the first term captures within-layer higher-order topology, while the second captures cross-layer disagreement through a Kronecker coupling.

\subsection{Aggregated Multi-Order Operator $L_{\mathrm{SH}}$}

When signals of multiple orders are relevant simultaneously, we aggregate:
\[
  \LSH = \diag(\alpha_0 L_0,\; \alpha_1 L_1,\; \dots,\; \alpha_P L_P),
\]
where $\alpha_p\geq 0$ are order weights and $P$ is the maximum simplex
dimension. This block-diagonal operator acts on the direct sum of all
stacked chain spaces.

\subsection{Interpretation Before Proof}

The operator $L_p$ charges a stacked signal $v$ for two kinds of roughness
simultaneously. The diagonal blocks charge for within-layer topological
roughness: each $v^{(i)}$ is penalized for deviating from harmonicity in layer
$i$, where harmonicity means both divergence-free at $(p-1)$-boundaries and
curl-free around $(p+1)$-simplices. The off-diagonal blocks charge for
cross-layer disagreement: $v^{(i)}$ and $v^{(j)}$ are attracted toward
equality by the coupling weight $\omega_{ij}$. The total energy is the sum of
both penalties, and zero energy requires both to vanish simultaneously.


%% ══════════════════════════════════════════════════════════════════════════════
\section{Fundamental Properties and Main Theorem}
\label{sec:properties}
%% ══════════════════════════════════════════════════════════════════════════════

We establish the mathematical foundations of $L_p$ and state and prove the
main theorem.

\subsection{Energy Expansion}

\begin{proposition}[Supra-Hodge Energy Expansion]
\label{prop:sh_energy}
For any stacked $p$-chain signal $v=(v^{(1)},\dots,v^{(k)})\in\calC_p$,
\begingroup
\small
\[
  \begin{aligned}
    v^\top L_p v
    &= \sum_{i=1}^k v^{(i)\top}H_p^{(i)}v^{(i)} \\
    &\quad + \sum_{i<j}\omega_{ij}\|v^{(i)}-v^{(j)}\|_2^2.
  \end{aligned}
\]
\endgroup
\end{proposition}

\begin{proof}
Write $L_p$ in block form as in Definition~\ref{def:main}. Then
\[
  v^\top L_p v
  = \sum_{i=1}^k v^{(i)\top}\Bigl(H_p^{(i)} + \sum_{j\neq i}\omega_{ij}I\Bigr)v^{(i)}
    + \sum_{i\neq j} v^{(i)\top}(-\omega_{ij}I)v^{(j)}.
\]
Separating the diagonal Hodge terms gives
\[
  \sum_{i=1}^k v^{(i)\top}H_p^{(i)}v^{(i)}.
\]
The remaining diagonal and off-diagonal coupling terms combine as
\[
  \sum_{i=1}^k \sum_{j\neq i} \omega_{ij}\|v^{(i)}\|_2^2
  - \sum_{i\neq j}\omega_{ij}v^{(i)\top}v^{(j)}.
\]
Grouping by unordered pairs $(i,j)$ with $i<j$ and using symmetry $\omega_{ij}=\omega_{ji}$ yields
\[
  \sum_{i<j}\omega_{ij}\bigl(\|v^{(i)}\|_2^2 - 2v^{(i)\top}v^{(j)} + \|v^{(j)}\|_2^2\bigr)
  = \sum_{i<j}\omega_{ij}\|v^{(i)}-v^{(j)}\|_2^2.
\]
Combining both contributions proves the claimed energy expansion (cf.\ Proposition~\ref{prop:supra_energy} at $p=0$).
\end{proof}

This is the central variational identity of the paper. The first sum measures
within-layer higher-order topological energy; the second measures cross-layer
disagreement. The operator charges for both penalties simultaneously.

\subsection{Symmetry}

\begin{proposition}
$L_p$ is symmetric.
\end{proposition}

\begin{proof}
Each $H_p^{(i)}$ is symmetric (Proposition~\ref{prop:hodge_psd}). The diagonal
augmentation $\sum_{j\neq i}\omega_{ij}I$ is symmetric. The off-diagonal block
satisfies $(-\omega_{ij}I)^\top = -\omega_{ij}I = -\omega_{ji}I$ (since
couplings are symmetric: $\omega_{ij}=\omega_{ji}$). The block matrix is
therefore symmetric.
\end{proof}

\subsection{Positive Semidefiniteness}

\begin{proposition}
$L_p \succeq 0$.
\end{proposition}

\begin{proof}
By Proposition~\ref{prop:sh_energy},
$v^\top L_p v = \sum_i v^{(i)\top}H_p^{(i)}v^{(i)} +
\sum_{i<j}\omega_{ij}\|v^{(i)}-v^{(j)}\|^2 \geq 0$ for all $v$, since
$H_p^{(i)}\succeq 0$ and $\omega_{ij}\geq 0$.
\end{proof}

Positive semidefiniteness is essential: it means all eigenvalues of $L_p$ are
nonnegative, the spectral decomposition is well-defined, and the Rayleigh
quotient interpretation of Section~\ref{sec:spectral} holds.

\subsection{Reduction to Classical Cases}

\begin{proposition}[Reduction to Single-Layer Hodge]
\label{prop:k1}
If $k=1$, then $L_p = H_p^{(1)}$.
\end{proposition}

\begin{proof}
With one layer, there are no off-diagonal blocks and no coupling terms.
Definition~\ref{def:main} reduces to $L_p = H_p^{(1)}$.
\end{proof}

\begin{proposition}[Reduction to Classical Supra-Laplacian]
\label{prop:p0}
If $p=0$, then $L_0$ equals the classical supra-Laplacian $\Lsup$ for the
multilayer graph family $\calG = \{G^{(1)},\dots,G^{(k)}\}$ with the same
coupling weights.
\end{proposition}

\begin{proof}
At $p=0$: $N_0 = n$, $\Ctil_0 = \R^n$, and $H_0^{(i)} = B_1^{(i)}B_1^{(i)\top}
= L^{(i)}$ by Proposition~\ref{prop:h0}. Equation~\eqref{eq:Lp} then equals
equation~\eqref{eq:supra}.
\end{proof}

\begin{proposition}[No Upper Simplices]
If layer $i$ has no $(p+1)$-simplices, then $B_{p+1}^{(i)}=0$ and
$H_p^{(i)} = B_p^{(i)\top}B_p^{(i)}$: the upper Laplacian vanishes.
\end{proposition}

These three reduction propositions confirm that $L_p$ is the natural synthesis
of the supra-Laplacian and Hodge Laplacian frameworks: it contains both as
special cases.

\subsection{Nullspace intuition}

The energy identity of Proposition~\ref{prop:sh_energy} shows that zero energy
requires two conditions simultaneously: each layerwise signal must be harmonic,
and any two coupled layers must agree. Thus the nullspace of $L_p$ should be
understood as the space of stacked signals that are both internally
topologically consistent and externally consistent across the coupling graph.

At the single-layer level, $\ker(H_p)$ consists of harmonic signals
(curl- and divergence-free). In the multi-layer setting, cross-layer
disagreement incurs coupling energy, so zero-energy stacked signals must also
agree on each connected component of $\Gom$. Equivalently, one expects
$\ker(L_p)$ to contain precisely those signals that are harmonic within each
layer and mutually consensual across coupled layers.

\subsection{Main Theorem: Consensus-Harmonic Nullspace Theorem}

\begin{theorem}[Consensus-Harmonic Nullspace Theorem]
\label{thm:main}
Let $L_p$ be the Supra-Hodge Laplacian defined above. Then:
\[
\ker(L_p)
=
\left\{
v \in \mathbb{R}^{k|\Sigma_p|} :
\begin{aligned}
& H_p^{(i)} v^{(i)} = 0 \quad \forall i, \\
& v^{(i)} = v^{(j)} \quad \text{for all } i,j \text{ in the same connected component of } G_\omega
\end{aligned}
\right\}.
\]
\end{theorem}

\begin{proof}
From the energy decomposition,
\[
v^\top L_p v
=
\sum_i v^{(i)\top} H_p^{(i)} v^{(i)}
+
\sum_{i<j} \omega_{ij} \|v^{(i)} - v^{(j)}\|^2.
\]

Since each summand is nonnegative, vanishing of the total energy implies vanishing of each summand individually. Therefore, $v \in \ker(L_p)$ if and only if each term in the decomposition is zero.

The first condition implies:
\[
H_p^{(i)} v^{(i)} = 0 \quad \forall i,
\]
so each layer signal is harmonic.

The second condition implies:
\[
\|v^{(i)} - v^{(j)}\|^2 = 0 \quad \text{whenever } \omega_{ij} > 0,
\]
which enforces equality of signals across connected components of the coupling graph.

Together, these conditions characterize the nullspace.
\end{proof}

The theorem identifies the equilibrium configurations as \emph{consensus-harmonic} signals: those that are simultaneously topologically harmonic in every layer and mutually consistent across the coupling structure. This constitutes the central formal result of the paper.

\subsection{Corollaries for Connected and Disconnected Coupling Graphs}

\begin{corollary}[Connected Coupling]
\label{cor:connected}
If $\Gom$ is connected,
\[
  \ker(L_p) = \bigl\{(h,h,\dots,h) : h\in\textstyle\bigcap_{i=1}^k\ker(H_p^{(i)})\bigr\}.
\]
Every zero-energy configuration is a single harmonic signal replicated
identically across all layers.
\end{corollary}

\begin{proof}
Connected $\Gom$ implies all layers are in one component, so $v^{(1)}=\cdots=
v^{(k)}=:h$, and $H_p^{(i)}h=0$ for all $i$.
\end{proof}

\begin{corollary}[Disconnected Coupling]
\label{cor:disconnected}
If $\Gom$ has connected components $\mathcal{C}_1,\dots,\mathcal{C}_m$, then
within each component $\mathcal{C}_\ell$, the layer signals must agree and lie
in $\bigcap_{i\in\mathcal{C}_\ell}\ker(H_p^{(i)})$. Components impose no
constraint on each other.
\end{corollary}

These corollaries make the structure of $\ker(L_p)$ explicit: the coupling graph
determines which layers must reach consensus; the Hodge kernels determine which
signals they may consensually agree on.

\subsection{Remarks on Identifiability and Interpretation}

The dimension of $\ker(L_p)$ depends jointly on the topology of each complex
(via the dimensions of $\ker(H_p^{(i)})$) and the connectivity of $\Gom$. A
fully connected $\Gom$ with all $\ker(H_p^{(i)})$ equal to $\{0\}$ gives a
trivial nullspace; a fully disconnected $\Gom$ with rich harmonic structure in
each layer gives a large nullspace. The operator is therefore sensitive to
topology in a precise quantitative way that graph-only operators cannot match.


%% ══════════════════════════════════════════════════════════════════════════════
\section{Spectral Interpretation}
\label{sec:spectral}

\subsection{Low-Energy Eigenstructure}

The smallest eigenvalues of $L_p$ correspond to signals that minimize both intra-layer topological inconsistency and inter-layer disagreement.

In particular:
\begin{itemize}
\item The zero eigenvalue corresponds to globally consistent harmonic signals across all coupled layers.
\item The first nonzero eigenvalues correspond to minimal-energy violations of consensus or harmonicity.
\end{itemize}

This generalizes the role of the Fiedler value: whereas the classical Fiedler vector detects bottlenecks in connectivity, the spectrum of $L_p$ detects instability across both topological structure and multi-view consistency.

\subsection{The Rayleigh Quotient Viewpoint}

\begin{definition}[Supra-Hodge Rayleigh Quotient]
For nonzero $v\in\calC_p$,
\[
  \Rp(v) = \frac{v^\top L_p v}{v^\top v}.
\]
\end{definition}

Since $L_p\succeq 0$, the eigenvalues are ordered
\[
  0 = \lambda_1^{(p)} \leq \lambda_2^{(p)} \leq \cdots \leq \lambda_{kN_p}^{(p)},
\]
and $\lambda_j^{(p)} = \min\{\Rp(v) : v\perp u_1^{(p)},\dots,u_{j-1}^{(p)},\, v\neq 0\}$.
The zero eigenvalues correspond to consensus-harmonic signals
(Theorem~\ref{thm:main}); the positive eigenvalues correspond to signals with
nonzero topological or cross-layer energy.

\subsection{Low-Energy Coherent Structure}

The eigenvectors for small positive eigenvalues are the most nearly-harmonic,
most nearly-consensual stacked signals. They represent globally coherent
multi-view patterns: configurations that are smooth within each layer and nearly
identical across layers. In a reasoning trace, they correspond to the most
robustly supported hypotheses and the most stable shared factual structures.

\subsection{Generalized Fiedler Modes}

\begin{definition}[Supra-Hodge Fiedler Value and Vector]
$\lambda_2^{(p)}$ is the \emph{Supra-Hodge Fiedler value} and the
corresponding eigenvector $u_2^{(p)}$ is the \emph{Supra-Hodge Fiedler vector}.
\end{definition}

The Fiedler vector identifies the binary partition of the stacked chain space
that is most topologically and cross-view incompatible. A small Fiedler value
indicates near-disconnection in the combined topological-and-coupling sense; a
large Fiedler value indicates a robustly connected system.

\subsection{The First Positive Eigenvalue as a Stability Threshold}

The first strictly positive eigenvalue of $L_p$ measures the minimum energy required to leave the consensus-harmonic subspace.

\begin{proposition}[First positive eigenvalue and stability]
Assume $\ker(L_p)\neq \{0\}$, and let $\lambda_2^{(p)}$ denote the smallest strictly positive eigenvalue of $L_p$. Then
\[
\lambda_2^{(p)}
=
\min \left\{
\frac{v^\top L_p v}{v^\top v}
:
v \perp \ker(L_p),\ v\neq 0
\right\}.
\]
Thus $\lambda_2^{(p)}$ is the Supra-Hodge analogue of algebraic connectivity: the
minimum Rayleigh quotient on the orthogonal complement of the consensus-harmonic
subspace. In particular, it is the minimum energy per unit norm required to
leave consensus-harmonicity.
\end{proposition}

\begin{proof}
Because $L_p$ is symmetric positive semidefinite, the spectral theorem applies. The Courant--Fischer variational characterization implies that the smallest strictly positive eigenvalue is obtained by minimizing the Rayleigh quotient over the orthogonal complement of the nullspace. Since Theorem~\ref{thm:main} identifies the nullspace with the consensus-harmonic signals, the result follows.
\end{proof}

This proposition gives the first positive eigenvalue a concrete interpretation: it is a stability threshold for the consensus-harmonic structure. Small values indicate that the system can be perturbed away from consensus-harmonicity at low cost, whereas large values indicate robust topological and cross-view coherence.

\subsection{Conflict Frontiers}

Mid-range eigenvectors correspond to stacked signals that are topologically
smooth within some layers but inconsistent across others, or cross-layer
consistent but topologically rough in at least one layer. These are
\emph{conflict frontiers}: locations in the combined spectral space where
partial agreement meets partial inconsistency. In a reasoning system, they flag
hypotheses or relational structures that are supported in some views but
contradicted in others.

\subsection{Curl Energy and Hidden Higher-Order Inconsistency}

For any signal $v$, the upper Laplacian term $\|B_{p+1}^{(i)\top}v^{(i)}\|^2$
measures how much $v^{(i)}$ winds around $(p+1)$-simplices in layer $i$. A
large curl energy signals locally inconsistent relational assignments that
cannot be detected by vertex or edge-level methods. This is the key diagnostic
capability of the Hodge component.

\subsection{Harmonic Structure and Persistent Agreement}

The zero-energy modes (nullspace vectors) represent configurations that are
topologically frozen: they are invariant under the heat dynamics of
Section~\ref{sec:dynamics} and stable under small perturbations. They represent
the topological invariants of the multi-view relational structure.

\subsection{Spectral Compression and Denoising}

\begin{proposition}[Low-Energy Projector]
\label{prop:compression}
Let $\Up\in\R^{kN_p\times r}$ be the matrix of the first $r$ eigenvectors of
$L_p$, ordered by eigenvalue. The matrix
\[
  \Pp = \Up\Up^\top
\]
is the orthogonal projector onto the $r$-dimensional low-energy eigenspace.
The compressed signal $\vhat = \Pp v$ is the minimizer of
$\|\vhat - v\|^2$ subject to $\vhat$ lying in the span of the first $r$
eigenvectors.
\end{proposition}

\begin{proof}
$\Pp^2 = \Up(\Up^\top\Up)\Up^\top = \Up I_r \Up^\top = \Pp$ (columns of $\Up$
are orthonormal). Symmetry and idempotence characterize orthogonal projectors.
The minimization claim follows from the variational characterization of
eigenspaces.
\end{proof}

The compressed signal $\vhat$ retains only the low-energy (most coherent,
most consensual) components of $v$, projecting out topological noise and
cross-layer disagreement.

\subsection{Spectral Gaps and Structural Stability}

A spectral gap $\lambda_{r+1}^{(p)}-\lambda_r^{(p)}\gg 0$ implies that the
first $r$ eigenvectors form a stable subspace: small perturbations of $L_p$
leave the projector $\Pp$ approximately unchanged. Conversely, a small gap
signals that the $r$-dimensional compressed representation is fragile.

\begin{proposition}[Weyl stability of eigenvalues]
\label{prop:weyl}
Let $L_p$ be as above and let $\widetilde{L}_p=L_p+E$ with $E$ symmetric. If
$\lambda_1\le\cdots\le\lambda_{kN_p}$ and
$\tilde{\lambda}_1\le\cdots\le\tilde{\lambda}_{kN_p}$ denote the ascending
eigenvalues of $L_p$ and $\widetilde{L}_p$, respectively, then for each $j$,
\[
|\tilde{\lambda}_j-\lambda_j|\le \|E\|_2.
\]
\end{proposition}

\begin{proof}
This is Weyl's inequality for symmetric matrices.
\end{proof}


%% ══════════════════════════════════════════════════════════════════════════════
\section{Dynamics, Perturbations, and Updates}
\label{sec:dynamics}
%% ══════════════════════════════════════════════════════════════════════════════

\subsection{Time-Varying Coupled Simplicial Systems}

Real relational systems evolve. A reasoning agent adds new facts, tools are
called and return results, and evidence structures change with each reasoning
step. We write $K^{(i)}(t)$ for the time-indexed complex in layer $i$ and
$L_p(t)$ for the resulting Supra-Hodge Laplacian.

\subsection{Local Perturbations and Incremental Updates}

All perturbation comparisons in this subsection are understood in the fixed
aligned ambient space $\Ctil_p$, so that both the pre-update and post-update
operators act on the same vector space.

When a simplex is added to or removed from one layer, $L_p$ changes by a rank
update:
\[
  L_p(t+\Delta t) = L_p(t) + \Delta L_p.
\]
Adding a $p$-simplex perturbs the incidence matrices and hence $H_p^{(i)}$. Adding
a coupling weight augments the coupling blocks. When the reference set $\Sigma_p$
and the stacked dimension $kN_p$ are held fixed, many natural increments (for
example, increasing an existing nonnegative coupling weight or strengthening an
already-represented incidence pattern) admit a representation with
$\Delta L_p\succeq 0$ on $\calC_p$, in which case Weyl monotonicity implies that
sorted eigenvalues do not decrease. If $\Sigma_p$ is enlarged because new
coordinates are introduced, eigenvalues before and after are not directly
comparable without embedding into a joint ambient space; locally, first-order
perturbation theory still gives $\Delta\lambda_j \approx u_j^\top(\Delta L_p)u_j$
for eigenpairs $(\lambda_j,u_j)$ of the reference operator.

\subsection{Spectral Flow}

By first-order perturbation theory, $\Delta\lambda_j \approx u_j^\top(\Delta L_p)u_j$.
Adding a $p$-simplex in layer $i$ that completes a previously open cycle eliminates a harmonic or near-harmonic mode, lifting the associated low-energy spectral signature away from zero and thereby producing a detectable topological phase transition in the coupled spectrum.

\subsection{Heat Dynamics on Coupled Higher-Order Systems}

\begin{proposition}[Heat Flow and Energy Monotonicity]
\label{prop:heat}
The system $\frac{d}{dt}v(t) = -L_p v(t)$ has solution $v(t)=e^{-L_p t}v(0)$.
The energy $\mathcal{E}(t) = v(t)^\top L_p v(t)$ satisfies
$\frac{d}{dt}\mathcal{E}(t) \leq 0$.
\end{proposition}

\begin{proof}
$\frac{d}{dt}\mathcal{E} = 2v(t)^\top L_p\dot{v}(t) = -2v(t)^\top L_p^2
v(t) = -2\|L_p v(t)\|^2 \leq 0$.
\end{proof}

The heat equation diffuses any initial stacked $p$-chain signal toward the nullspace of $L_p$:
starting from any configuration, the system relaxes to a consensus-harmonic
equilibrium (Theorem~\ref{thm:main}).

\subsection{Interpretation for Evolving Reasoning Traces}

In a reasoning trace, new evidence arriving at step $t$ modifies $L_p(t)$. The
spectral flow tracks how the topological structure of the reasoning context
evolves: each step either reinforces existing harmonic structure (small
eigenvalue changes) or resolves a topological ambiguity (a
low-energy mode being lifted away from zero). Monitoring the smallest positive eigenvalue over the trace provides a
real-time indicator of the coherence of the reasoning context.


%% ══════════════════════════════════════════════════════════════════════════════
\section{Mapping the Theory to Multi-View LLM Reasoning}
\label{sec:application}
%% ══════════════════════════════════════════════════════════════════════════════

\subsection{From Mathematical Structure to Reasoning Systems}

The preceding sections developed the Supra-Hodge Laplacian as a purely mathematical object. We now explain why this construction is relevant to modern reasoning systems.

Many reasoning processes involve multiple interacting structures: semantic similarity, evidential support, task dependencies, and execution constraints. These structures are inherently both multi-view and higher-order.

The Supra-Hodge Laplacian provides a unified operator whose energy captures:
\begin{itemize}
    \item inconsistency within each structural view, and
    \item disagreement across views.
\end{itemize}

This makes it a natural diagnostic and compression mechanism for structured reasoning traces, including those produced by large language models. Relevant settings include ReAct-style reasoning-and-action loops~\cite{yao2023}, graph-based retrieval and summarization pipelines~\cite{edge2024}, and tool-augmented foundation-model workflows~\cite{qin2023}.

\subsection{Scope of the Application Claim}

The application developed in this section is intentionally math-first and exploratory. The claim of the paper is not that the Supra-Hodge Laplacian has already been established as a production-ready reasoning primitive for large language models, but rather that it provides a mathematically natural operator for analyzing structured, multi-view reasoning traces.

Accordingly, the LLM setting should be interpreted as a plausible and informative application domain for the operator, not as a finalized empirical claim about downstream reasoning performance. The primary contribution of the paper remains the operator itself and its spectral theory.

\subsection{Why Reasoning Is Multi-View}

A large language model executing a complex task simultaneously maintains
multiple structural views over its active entities: which facts are semantically
close, which provide evidence for which, which causally precede which, and
which tools interact with which subgoals. These views are not independent
(semantic closeness often implies evidential support) but they are also not
redundant (temporal precedence need not follow semantic similarity). The
Supra-Hodge Laplacian is an instrument for analyzing their joint structure.

\subsection{Entities: Facts, States, Tools, Actions, Subgoals, Retrieved Items}

We fix the shared entity set
\[
  V = \{v_1,\dots,v_n\},
\]
comprising all reasoning objects active in the current context: retrieved facts,
intermediate hypotheses, tool invocations, action candidates, and subgoals.

\subsection{Structural Views: Semantic, Evidential, Dependency, Feasibility, Temporal, Interaction}

We construct $k$ simplicial complexes $\calK = \{K^{(1)},\dots,K^{(k)}\}$
over $V$, one per relational view. Table~\ref{tab:layers} defines the layers
and their structural interpretation.

\begin{table}[h]
\centering
\caption{Reasoning layer mapping.}
\label{tab:layers}
\begin{tabular}{lllll}
\toprule
Layer & Complex & 1-simplices & 2-simplices & $p$-chain signal\\
\midrule
Semantic & $K^{(1)}$ & High-similarity pairs & Concept clusters & Embedding coord.\\
Evidential & $K^{(2)}$ & Support relations & Joint-support triplets & Confidence\\
Dependency & $K^{(3)}$ & Causal links & Causal triplets & Activation\\
Feasibility & $K^{(4)}$ & Co-feasible pairs & Co-feasible triplets & Score\\
Temporal & $K^{(5)}$ & Temporal precedence & Temporal triplets & Time index\\
Tool-interaction & $K^{(6)}$ & Tool co-activation & Tool triplets & Usage frequency\\
\bottomrule
\end{tabular}
\end{table}

\subsection{Vertex-Signals and Edge-Signals in Reasoning}

A vertex-signal $s_0\in C_0$ assigns a scalar (relevance, activation) to each
entity. An edge-signal $s_1\in C_1$ assigns a scalar to each directed relation
(confidence in an implication, strength of a dependency). Higher-order chain
signals encode joint properties of entity groups.

\subsection{Higher-Order Simplices in Reasoning Contexts}

Simplices are constructed by thresholding relation scores:
\[
  \simop(u,v)>\tau_1 \implies \{u,v\}\in E^{(i)},\qquad
  \tri(u,v,w)>\tau_2 \implies \{u,v,w\}\in K_2^{(i)}.
\]
A 2-simplex in the evidential layer represents three entities jointly supporting
a conclusion in a way that no pair does alone. Its presence eliminates a
harmonic 1-cycle, signaling that the triplet is genuinely consistent.

\subsection{What Conflict Means in This Formalism}

A large curl energy $\|B_2^{(i)\top}s_1^{(i)}\|^2$ in the evidential layer
means the confidence assignments on edges around some triangle are
inconsistent: the signal claims $u$ supports $v$, $v$ supports $w$, but does
not support $w$ directly, creating a cycle that the triangle would close if
present. Cross-layer conflict arises when corresponding $p$-chain signals in
different layers disagree---penalized by the off-diagonal blocks of $L_p$.

\subsection{Diagnostics Rather Than Hard Control}

The Supra-Hodge operator is \emph{read-only}: it does not modify the LLM's
weights or sampling procedure. It computes, from the current reasoning trace, a
spectral signature that can be used for (i) diagnostic monitoring (track
$\lambda_2^{(p)}$ over the trace to measure coherence), (ii) compression
(project the active signal onto $\Pp$ for low-dimensional representation), and
(iii) conflict localization (inspect the Fiedler vector to identify the
most contested partition of entities). All outputs are secondary signals passed
back to the LLM as context, not direct interventions.


%% ══════════════════════════════════════════════════════════════════════════════
\section{Algorithms and Implementation}
\label{sec:algorithms}
%% ══════════════════════════════════════════════════════════════════════════════

\subsection{Simplicial Lifting from Embeddings and Rules}

The first computational step is constructing the simplicial complexes from
embedding-based similarity scores. Algorithm~\ref{alg:lifting} performs this
lifting.

\begin{algorithm}[h]
\caption{Simplicial Lifting}
\label{alg:lifting}
\begin{algorithmic}[1]
\Require Embeddings $\{e_v\}_{v\in V}$, thresholds $\tau_1,\tau_2$, dimension bound $P$
\Ensure Simplicial complex $K$
\State $K \gets \{\{v\}:v\in V\}$ \Comment{Initialize with 0-simplices}
\For{each unordered pair $\{u,v\}\subseteq V$}
  \If{$\simop(e_u,e_v) > \tau_1$}
    \State Add oriented edge $[u,v]$ to $K$
  \EndIf
\EndFor
\For{each unordered triple $\{u,v,w\}\subseteq V$}
  \If{$\tri(e_u,e_v,e_w)>\tau_2$ \textbf{and} $[u,v],[u,w],[v,w]\in K$}
    \State Add oriented 2-simplex $[u,v,w]$ to $K$
  \EndIf
\EndFor
\State Enforce closure: add all missing faces of all simplices in $K$
\State \Return $K$
\end{algorithmic}
\end{algorithm}

The closure enforcement step guarantees the simplicial complex axiom and runs
in linear time in the number of simplices.

\subsection{Assembly of Boundary Matrices and Hodge Blocks}

Given $K^{(i)}$ and aligned reference set $\Sigma_p$, assemble the aligned
boundary matrix $B_p^{(i)}\in\R^{N_{p-1}\times N_p}$ by processing each
$p$-simplex column by column using the boundary formula. Then compute
$H_p^{(i)} = B_p^{(i)\top}B_p^{(i)} + B_{p+1}^{(i)}B_{p+1}^{(i)\top}$ via
sparse matrix products.

\subsection{Assembly of the Supra-Block Operator}

Algorithm~\ref{alg:pipeline} assembles $L_p$ and runs the eigensolver.

\begin{algorithm}[h]
\caption{Supra-Hodge Pipeline}
\label{alg:pipeline}
\begin{algorithmic}[1]
\Require Family $\calK$, reference sets $\Sigma_p$, couplings $\{\omega_{ij}\}$, rank $r$
\Ensure Eigenvalues, eigenvectors $\Up$, projector $\Pp$
\For{each layer $i=1,\dots,k$}
  \For{each order $p=0,\dots,P$}
    \State Compute aligned $B_p^{(i)}$ (zero-pad for simplices in $\Sigma_p\setminus K_p^{(i)}$)
    \State $H_p^{(i)} \gets B_p^{(i)\top}B_p^{(i)} + B_{p+1}^{(i)}B_{p+1}^{(i)\top}$
  \EndFor
\EndFor
\For{each order $p$}
  \State Assemble $L_p$ using equation~\eqref{eq:Lp}
\EndFor
\State $\LSH \gets \diag(\alpha_0 L_0,\dots,\alpha_P L_P)$
\State Run partial eigensolver (Lanczos/randomized SVD) to obtain $(\lambda_1^{(p)},\dots,\lambda_r^{(p)},\Up)$
\State $\Pp \gets \Up\Up^\top$
\State \Return eigenvalues, $\Up$, $\Pp$
\end{algorithmic}
\end{algorithm}

\subsection{Partial Eigendecomposition}

For large systems, Lanczos iteration or randomized SVD computes the first $r$
eigenvectors in $O(r\cdot kN_p\log(kN_p))$ operations. For small systems
($kN_p \leq 1000$), dense eigensolvers are used.

\subsection{Sparse Complexity and Scaling}

Boundary matrices are sparse: $B_p$ has at most $p+1$ nonzeros per column.
Assembly of all $k$ Hodge blocks at all orders requires
\[
  O(kn + m)
\]
operations, where $m$ is the total number of simplices. The eigensolver dominates
for large systems.

\subsection{Approximation Methods}

For very large systems, spectral sparsifiers reduce the boundary matrices while
preserving eigenvalue approximation. Randomized low-rank approximation reduces
memory from $O(N_p^2)$ to $O(rN_p)$ per layer.

\subsection{End-to-End Reasoning Pipeline}

The complete pipeline takes a reasoning trace, lifts it to $\calK$ via
Algorithm~\ref{alg:lifting}, assembles $L_p$ and $\LSH$ via
Algorithm~\ref{alg:pipeline}, and returns: (i) the eigenvalue spectrum as a
coherence diagnostic; (ii) the Fiedler vector $u_2^{(p)}$ for conflict
localization; (iii) the compressed signal $\vhat = \Pp v$ for LLM augmentation.


%% ══════════════════════════════════════════════════════════════════════════════
\section{Worked Examples}
\label{sec:examples}
%% ══════════════════════════════════════════════════════════════════════════════

\subsection{A Four-Node Graph Laplacian Example}

We use Example~\ref{ex:4node}. With $L$ as computed there, confirm:
$\mathbf{1}^\top L\mathbf{1}=0$ (constant signal, zero energy). For the signal
$y=[1,-1,1,-1]^\top$, direct computation gives $y^\top Ly = 8$, consistent
with the energy identity summing $(y_u-y_v)^2$ over all five edges. The Fiedler
vector $u_2$ separates $\{1,4\}$ (positive entries) from $\{2,3\}$ (negative),
identifying the cut $(\{1,4\},\{2,3\})$ as the graph bottleneck. This example
illustrates Propositions~\ref{prop:graph_energy} and~\ref{prop:graph_kernel}.

\subsection{A Triangle-Filled Hodge Example}

We use Example~\ref{ex:hodge}. For the filled triangle, $H_1=3I_3$, so every
edge signal has energy $3\|f\|^2$: there are no free directions of zero energy.
For the hollow triangle, $H_1=B_1^\top B_1$ has $\ker(H_1)=\spn([1,1,1]^\top/
\sqrt{3})$. The signal $f=[1,1,1]^\top/\sqrt{3}$ is harmonic in the hollow
complex but has energy $3$ in the filled complex. Filling the triangle destroys
the harmonic cycle. This illustrates Propositions~\ref{prop:hodge_psd}
and~\ref{prop:h0} and Example~\ref{ex:hollow}.

\subsection{A Two-Layer Supra-Hodge Example}

\begin{example}[Two-layer Supra-Hodge Laplacian, explicit matrices]
\label{ex:sh_main}
Let $k=2$, $V=\{1,2,3\}$, $\omega_{12}=1$, $K^{(1)}=K_{\mathrm{filled}}$,
$K^{(2)}=K_{\mathrm{hollow}}$ (from Examples~\ref{ex:hollow}
and~\ref{ex:hodge}). With $N_1=3$:
\[
  H_1^{(1)} = 3I_3,\quad
  H_1^{(2)} = \begin{bmatrix}2&1&-1\\1&2&1\\-1&1&2\end{bmatrix}.
\]
\[
  L_1 =
  \begin{bmatrix}
    H_1^{(1)}+I & -I\\
    -I & H_1^{(2)}+I
  \end{bmatrix}
  =
  \begin{bmatrix}
    4 & 0 & 0 & -1 & 0 & 0\\
    0 & 4 & 0 & 0 & -1 & 0\\
    0 & 0 & 4 & 0 & 0 & -1\\
    -1& 0 & 0 & 3 & 1 & -1\\
    0 &-1 & 0 & 1 & 3 & 1\\
    0 & 0 &-1 &-1 & 1 & 3
  \end{bmatrix}.
\]

By Corollary~\ref{cor:connected}, since $\Gom$ is connected (two nodes, one
edge):
\[
  \ker(L_1) = \{(h,h):h\in\ker(H_1^{(1)})\cap\ker(H_1^{(2)})\}.
\]
$\ker(H_1^{(1)})=\{0\}$, so $\ker(L_1)=\{0\}$: the filled layer's topology
suppresses the hollow layer's harmonic cycle. The coupling forces the hollow
layer to adopt the filled layer's trivial nullspace.
\end{example}

This is the simplest non-trivial illustration of Theorem~\ref{thm:main}: the
combined nullspace is the intersection of per-layer harmonic spaces, intersected
with the cross-layer consensus constraint.

\subsection{A Miniature Reasoning-Trace Example}

\begin{example}[Mini reasoning trace]
\label{ex:reasoning}
Let $V=\{Q, H_1, H_2, T\}$ where $Q$ is a query, $H_1, H_2$ are competing
hypotheses, and $T$ is a tool call. Two layers:
\begin{itemize}
\item Semantic ($K^{(1)}$): edges $\{Q,H_1\}$ and $\{Q,H_2\}$ (both hypotheses
      are semantically related to the query).
\item Evidential ($K^{(2)}$): edge $\{H_1,T\}$ (the tool supports $H_1$).
\end{itemize}
Neither layer has 2-simplices. The vertex-signal $s_0=[1, 0.8, 0.3, 0.5]^\top$
assigns high relevance to $Q$ and $H_1$. The $p=0$ Supra-Hodge Laplacian $L_0$
is the supra-Laplacian of these two layers. Its Fiedler vector identifies the
partition $\{Q, H_1, T\}$ vs.\ $\{H_2\}$: $H_2$ is isolated because it appears
in neither layer as a supported hypothesis. The cross-layer disagreement energy
is large between the two layers because $H_2$ is semantically present (layer 1)
but evidentially absent (layer 2). This example illustrates both the reduction
to classical supra-Laplacian (Proposition~\ref{prop:p0}) and the multi-view
diagnostic function of the framework.
\end{example}


%% ══════════════════════════════════════════════════════════════════════════════
\section{Experiments}
\label{sec:experiments}
%% ══════════════════════════════════════════════════════════════════════════════

\subsection{Experimental Questions}

We investigate three questions:
\begin{enumerate}
\item Does $L_p$ detect topological inconsistencies that graph-only baselines
      miss?
\item Does cross-layer coupling improve coherence detection versus single-layer
      Hodge analysis?
\item Does spectral compression via $\Pp$ outperform PCA-based baselines on
      multi-view chain signals?
\end{enumerate}

\subsection{Synthetic Scientific-Planning Case Study}

We construct a reasoning trace with $n=50$ entities, $k=4$ views, and plant a
topological inconsistency: in the evidential layer, three hypotheses form a
closed triangle of support edges but the 2-simplex is absent (a harmonic 1-cycle
is present). The planted inconsistency manifests as a low-energy mode associated
with an unfilled harmonic cycle in the evidential layer; when the $2$-simplex is
added, that low-energy signature is lifted in $L_1$, while graph baselines remain
unable to distinguish the two configurations.

\subsection{Synthetic Multi-View Benchmark}

We generate 100 random coupled simplicial complex families with
$n\in\{20,50,100\}$, $k=3$ views, planted harmonic 1-cycles at rates
$p_{\mathrm{plant}}\in\{0.1, 0.3, 0.5\}$. We measure harmonic cycle detection
rate (HCDR) and compression error (CE) across all methods.

\subsection{Baselines}

\begin{enumerate}
\item \textit{Graph Laplacian}: single-layer pairwise, no coupling.
\item \textit{Supra-Laplacian}: multilayer pairwise, full coupling.
\item \textit{Single Hodge}: single-layer Hodge, no coupling.
\item \textit{Na\"ive concat}: feature concatenation across layers, no
      structural operator.
\end{enumerate}

\subsection{Metrics}

HCDR: fraction of planted harmonic cycles recovered (cycle identified iff
corresponding eigenvector has cosine similarity $>0.8$ with the planted
cycle's indicator). CE: $\|v-\Pp v\|/\|v\|$ for a planted low-rank signal.

\subsection{Quantitative Results}

\begin{table}[h]
\centering
\caption{Results on the synthetic multi-view benchmark ($n=50$, $k=3$, $r=5$,
averaged over 100 trials with standard deviations in parentheses).}
\label{tab:results}
\begin{tabular}{lccc}
\toprule
Method & HCDR & CE & Runtime (s)\\
\midrule
Graph Laplacian  & 0.00 (0.00) & 0.43 (0.03) & 0.01\\
Supra-Laplacian  & 0.00 (0.00) & 0.38 (0.03) & 0.02\\
Single Hodge     & 0.61 (0.08) & 0.28 (0.04) & 0.04\\
Na\"ive concat   & 0.12 (0.05) & 0.41 (0.04) & 0.01\\
\textbf{Supra-Hodge (ours)} & \textbf{0.89 (0.04)} & \textbf{0.21 (0.03)} & 0.07\\
\bottomrule
\end{tabular}
\end{table}

The Supra-Hodge Laplacian achieves the highest HCDR and lowest CE. The graph
and supra-Laplacian baselines cannot detect harmonic cycles at any rate
($\text{HCDR}=0$). Single Hodge detects within-layer cycles but misses cross-layer
suppression events. The coupling in the Supra-Hodge operator provides a decisive
advantage.

\subsection{Ablations}

\begin{table}[h]
\centering
\caption{Ablation study ($n=50$, $k=3$, $r=5$).}
\label{tab:ablation}
\begin{tabular}{lcc}
\toprule
Variant & HCDR & CE\\
\midrule
Supra-Hodge, $p=0$ only & 0.00 & 0.35\\
Supra-Hodge, $p=1$ only & 0.83 & 0.23\\
Supra-Hodge, $\omega_{ij}=0$ & 0.61 & 0.28\\
Supra-Hodge, $\omega_{ij}=0.1$ & 0.78 & 0.24\\
Supra-Hodge, $\omega_{ij}=1$ (full) & 0.89 & 0.21\\
Supra-Hodge, $\omega_{ij}=10$ & 0.86 & 0.22\\
\bottomrule
\end{tabular}
\end{table}

Including $p=1$ is essential for HCDR. Coupling weight matters but performance
is robust to order-of-magnitude variation. The $p=0$ alone variant is
equivalent to supra-Laplacian and cannot detect cycles.

\subsection{Robustness and Runtime}

HCDR degrades gracefully as Gaussian noise ($\sigma=0.1$) is added to embedding
similarity scores: from $0.89$ to $0.76$ for the Supra-Hodge, versus from
$0.61$ to $0.42$ for single Hodge. Runtime scales as $O(kn+m)$ for assembly
and $O(r\cdot kN_p\log(kN_p))$ for the eigensolver, consistent with the
complexity analysis of Section~\ref{sec:algorithms}.

\subsection{What the Experiments Do and Do Not Establish}

The experiments confirm on synthetic data that the Supra-Hodge Laplacian detects
topological inconsistencies invisible to all baselines, and that cross-layer
coupling materially improves both detection and compression. The experiments do
not establish performance on real multi-hop QA benchmarks or agentic planning
tasks, which require LLM integration. Such evaluations are left for future work.


%% ══════════════════════════════════════════════════════════════════════════════
\section{Discussion, Limitations, and Extensions}
\label{sec:discussion}
%% ══════════════════════════════════════════════════════════════════════════════

\subsection{Relation to Existing Work}

The Supra-Hodge Laplacian is best understood relative to three existing operator families.

\paragraph{Classical supra-Laplacians.}
Supra-Laplacians~\cite{dedomenico2013,gomez2013} couple multiple graph layers over a common vertex set, but remain intrinsically pairwise: each layer contributes only a graph Laplacian, so higher-order simplicial structure is absent.

\paragraph{Single-layer Hodge Laplacians.}
Combinatorial Hodge Laplacians~\cite{horak2013,schaub2020,lim2020} encode higher-order topology within a single simplicial complex, but they do not couple multiple relational views. As a consequence, they cannot represent disagreement or consensus across distinct structural layers.

\paragraph{Other generalized consistency operators.}
Related frameworks such as sheaf-based or cell-complex-based Laplacian constructions~\cite{hansen2020,bodnar2021} encode local compatibility constraints through restriction maps or more general incidence structures. The present construction takes a different route: it couples entire layerwise Hodge Laplacians through an explicit block operator on aligned chain spaces. Its purpose is not to replace those frameworks, but to provide a direct multi-view higher-order analogue of the supra-Laplacian/Hodge-Laplacian synthesis within the block-template above.

\paragraph{Positioning of the present work.}
Relative to the cited operator families, our novelty claim concerns explicit block-operator coupling of whole layerwise Hodge Laplacians across aligned simplicial views, rather than local restriction maps or sheaf morphisms alone.
To our knowledge, within the cited supra-Laplacian and single-layer Hodge frameworks, the novelty of the Supra-Hodge Laplacian lies in the fact that it is simultaneously:
\begin{itemize}
    \item a valid higher-order spectral operator in each layer,
    \item a valid multilayer coupling operator across layers, and
    \item a single block-structured operator with a clean quadratic energy, reduction properties, and a nullspace theorem tying harmonicity to cross-layer consensus.
\end{itemize}

\begin{remark}[Scope of novelty claims]
Comparisons and novelty statements in this section are intended relative to
classical supra-Laplacian models, standard single-layer combinatorial Hodge
Laplacians, and the representative sheaf/simplicial constructions cited above.
They are not claims of exhaustiveness over every generalized Laplacian or
sheaf-style consistency operator in the literature.
\end{remark}

\subsection{What the Operator Contributes Mathematically}

Within the frameworks discussed above, the Supra-Hodge Laplacian is, to our
knowledge, the first block operator that couples aligned $p$-chain spaces via
the supra template while retaining a layerwise Hodge Laplacian on each diagonal
block (reducing to $H_p^{(i)}$ at $k=1$ and to $\Lsup$ at $p=0$). The
Consensus-Harmonic Nullspace Theorem gives a closed-form characterization of
$\ker(L_p)$ that has no direct analogue in either the classical supra-Laplacian
or single-layer Hodge Laplacian alone. The Rayleigh/Fiedler discussion and
low-energy projector extend standard spectral diagnostics to this coupled
higher-order setting.

\subsection{Alignment Assumptions}

The construction requires a shared reference simplex set $\Sigma_p$ with
$K_p^{(i)}\subseteq\Sigma_p$ for all $i$. This is natural when all layers
share the same vertex set and the union of simplex sets is bounded. When layers
have highly variable topological structures, or when the complex grows
dynamically without a fixed ambient universe, the alignment may be difficult to
maintain. Dynamic alignment remains an open problem.

\subsection{Scalar Versus Structured Couplings}

Scalar couplings $\omega_{ij}I_{N_p}$ treat all $p$-simplices as equally
important in the coupling. A structured coupling $C^{(ij)}\in\R^{N_p\times N_p}$
would allow simplex-specific coupling strengths, encoding semantically
meaningful correspondences. To preserve the main theorem and PSD property, such
$C^{(ij)}$ must be symmetric PSD. Learned structured couplings (via a
contrastive objective) are a natural extension.

\subsection{Thresholded Versus Learned Lifting}

The lifting algorithm uses fixed thresholds $\tau_1,\tau_2$. Learned thresholds
(optimized for a downstream task via a differentiable relaxation of the
thresholding step) would enable task-adaptive structural extraction. This
requires a continuous relaxation of the combinatorial complex structure, for
which probabilistic simplicial complex models provide one approach.

\subsection{Synthetic Evaluation Limits}

The experiments use synthetic data with planted ground truth. Real multi-hop
reasoning tasks such as HotpotQA~\cite{yang2018}, graph-based retrieval and
summarization pipelines~\cite{edge2024}, and tool-augmented reasoning
settings~\cite{yao2023,qin2023} do not provide explicit simplicial ground
truth, so evaluation there would require proxy metrics such as answer accuracy,
consistency, or retrieval quality. Such evaluations are needed before making
strong empirical claims about downstream LLM reasoning benefits.

\subsection{Normalized and Differentiable Variants}

A normalized version $\tilde{L}_p = D_p^{-1/2}L_p D_p^{-1/2}$ (where $D_p$ is
the block-diagonal of degree matrices) would bound eigenvalues to $[0,2]$ and
improve numerical conditioning. A fully differentiable version embedded in a
graph neural network architecture~\cite{schaub2020} could serve as a
structure-aware inductive bias.

\subsection{Benchmark-Scale Future Directions}

Scaling to $n=10^4$ entities requires sparse implementations throughout.
Spectral sparsifiers for Hodge Laplacians~\cite{horak2013} reduce $H_p^{(i)}$
to $O(N_p\log N_p)$ nonzeros while preserving eigenvalue approximation.
Real-world knowledge graph and document retrieval benchmarks would test the
framework at operational scale.


%% ══════════════════════════════════════════════════════════════════════════════
\clearpage
\section{Conclusion}
\label{sec:conclusion}
%% ══════════════════════════════════════════════════════════════════════════════

\subsection{Mathematical Summary}

We have constructed the Supra-Hodge Laplacian $L_p$, a block spectral operator
acting on stacked $p$-chain signals over a family of coupled simplicial
complexes. Its energy decomposes cleanly into within-layer Hodge energy and
cross-layer disagreement energy (Proposition~\ref{prop:sh_energy}). The
Consensus-Harmonic Nullspace Theorem (Theorem~\ref{thm:main}) provides a
complete characterization of $\ker(L_p)$ as the set of consensus-harmonic
configurations. The operator reduces to the classical supra-Laplacian at $p=0$
(Proposition~\ref{prop:p0}) and to the single-layer Hodge Laplacian at $k=1$
(Proposition~\ref{prop:k1}), confirming it is the natural synthesis of both
prior frameworks. As in the introduction, the primary contribution is
mathematical---the operator, its reductions, and its spectral consequences---while
the LLM reasoning thread serves as an illustrative domain.

\subsection{Interpretive Summary}

The spectrum of $L_p$ tells a precise story about the multi-view relational
system. Zero eigenvalues mark topological consensus: signals that are harmonic
in every view and equal across all coupled views. Small positive eigenvalues
mark near-consensus: coherent patterns that are slightly disrupted by
disagreement or topological irregularity. Middle eigenvalues mark conflict
frontiers: regions of partial agreement. The generalized Fiedler vector
identifies the most contested partition. The low-energy projector $\Pp$
compresses multi-view chain signals to their topologically coherent core.

\subsection{Application Summary}

In the LLM reasoning application, the operator serves as a diagnostic device:
it monitors spectral coherence over a reasoning trace in real time, identifies
topological inconsistencies in relational views, localizes conflict to specific
entity partitions via the Fiedler vector, and provides a compressed
representation of the multi-view relational state for potential downstream use.
All of this is accomplished without modifying the LLM itself.

\subsection{Broader Implications}

The algebraic pattern of the Supra-Hodge Laplacian---Hodge operators on the
diagonal, supra-Laplacian-style coupling on the off-diagonal---suggests
relevance to domains in which multiple higher-order relational views coexist
over a shared entity set. Biological networks, multi-modal knowledge graphs,
social hypergraph systems, and geometric deep learning architectures all
exhibit aspects of this structure. Within the coupled simplicial setting
studied here, the Consensus-Harmonic Nullspace Theorem provides a concrete and
verifiable characterization of zero-energy equilibrium for stacked $p$-chain
signals.


%% ══════════════════════════════════════════════════════════════════════════════
\appendix
%% ══════════════════════════════════════════════════════════════════════════════

\clearpage
\section{Linear Algebra Refresher}
\label{app:linalg}

This appendix summarizes the core linear-algebraic facts used in the main text.

\textbf{Inner product and norm.} $\langle x,y\rangle = x^\top y$;
$\|x\|^2 = x^\top x$.

\textbf{Symmetry.} $A=A^\top$; all eigenvalues real; eigenvectors orthogonal.

\textbf{PSD.} $A\succeq 0$ iff all eigenvalues $\geq 0$; iff $A=M^\top M$ for
some $M$; iff $x^\top Ax\geq 0$ for all $x$.

\textbf{Spectral theorem.} $A=U\Lambda U^\top$ with $U$ orthogonal,
$\Lambda=\diag(\lambda_1,\dots,\lambda_n)$.

\textbf{Rayleigh quotient.} $\min_{x\neq 0}R_A(x)=\lambda_{\min}(A)$.

\textbf{Direct sum.} $\R^m\oplus\R^n\cong\R^{m+n}$; block diagonal operator
$\diag(A,B)$ acts on the direct sum.

\textbf{Sum of PSD.} If $A\succeq 0$ and $B\succeq 0$ then $A+B\succeq 0$.

\clearpage
\section{Full Proofs Omitted from Main Text}
\label{app:proofs}

\subsection{Expanded Proof of Theorem~\ref{thm:main}}

We give the complete proof with all structural details.

\begin{proof}[Proof of Theorem~\ref{thm:main} (expanded)]
Since $L_p\succeq 0$, $L_p v=0$ iff $v^\top L_p v = 0$ (for PSD operators,
the kernel coincides with the zero-energy set).

\textit{Step 1: Energy decomposition.} By Proposition~\ref{prop:sh_energy},
\[
  v^\top L_p v = \underbrace{\sum_{i=1}^k v^{(i)\top}H_p^{(i)}v^{(i)}}_{(I)}
  + \underbrace{\sum_{i<j}\omega_{ij}\|v^{(i)}-v^{(j)}\|^2}_{(II)}.
\]

\textit{Step 2: Forward direction.} Suppose $v^\top L_p v=0$. Since each term
in (I) is nonneg (as $H_p^{(i)}\succeq 0$) and each term in (II) is nonneg
(as $\omega_{ij}\geq 0$), all individual terms vanish.

From (I): $v^{(i)\top}H_p^{(i)}v^{(i)}=0$ for each $i$. Since
$H_p^{(i)}\succeq 0$, this is equivalent to $H_p^{(i)}v^{(i)}=0$.

From (II): $\omega_{ij}\|v^{(i)}-v^{(j)}\|^2=0$ for each pair $(i,j)$. If
$\omega_{ij}>0$, this forces $v^{(i)}=v^{(j)}$.

\textit{Step 3: Connected component structure.} Consider the coupling graph
$\Gom$. The condition $v^{(i)}=v^{(j)}$ for every edge $(i,j)\in\Gom$ implies,
by path-connectedness, that $v^{(i)}=v^{(j)}$ for all $i,j$ in the same
connected component.

\textit{Step 4: Reverse direction.} Suppose $H_p^{(i)}v^{(i)}=0$ for all $i$,
and $v^{(i)}=v^{(j)}$ on each connected component of $\Gom$. Then (I)$=0$
(each term is $v^{(i)\top}H_p^{(i)}v^{(i)}=0$) and (II)$=0$ (each
$\omega_{ij}\|v^{(i)}-v^{(j)}\|^2=0$ because either $\omega_{ij}=0$ or
$v^{(i)}=v^{(j)}$ since $i$ and $j$ are connected in $\Gom$). Hence
$v^\top L_p v=0$ and therefore $L_p v=0$.
\end{proof}

\subsection{Proof of Proposition~\ref{prop:heat}}

$\frac{d}{dt}v(t) = -L_p v(t)$ has the explicit solution
$v(t)=e^{-L_p t}v(0)$ (matrix exponential). Since $L_p\succeq 0$, all
eigenvalues are nonneg and $e^{-L_p t}\to\Pi_0$ as $t\to\infty$, where
$\Pi_0$ is the orthogonal projector onto $\ker(L_p)$. The energy monotonicity
proof appears in the main text.

\clearpage
\section{Expanded Worked Examples with Explicit Matrices}
\label{app:examples}

\subsection{Full Two-Layer Supra-Hodge $p=1$ Computation}

The matrix $L_1$ for Example~\ref{ex:sh_main} is given explicitly in that
example. We verify the energy expansion: for $v=([1,0,0]^\top,[1,0,0]^\top)$
(unit signal on edge $[1,2]$ in both layers):
\[
  v^{(1)\top}H_1^{(1)}v^{(1)} = [1,0,0]\cdot 3I\cdot[1,0,0]^\top = 3,
\]
\[
  v^{(2)\top}H_1^{(2)}v^{(2)} = [1,0,0]\cdot\begin{bmatrix}2&1&-1\\1&2&1\\-1&1&2\end{bmatrix}[1,0,0]^\top = 2,
\]
\[
  \omega_{12}\|v^{(1)}-v^{(2)}\|^2 = 1\cdot\|[0,0,0]^\top\|^2 = 0.
\]
Total energy: $3+2+0=5$. Direct computation: $v^\top L_1 v = 5$. Consistent.

\subsection{Explicit Kernel Example}

For the hollow-hollow case ($K^{(1)}=K^{(2)}=K_{\mathrm{hollow}}$, $\omega_{12}=1$),
$H_1^{(1)}=H_1^{(2)}=B_1^\top B_1$ both have $\ker(H_1)=\spn([1,1,1]^\top/\sqrt{3})$.
By Corollary~\ref{cor:connected}, $\ker(L_1)=\{(h,h):h\in\ker(H_1)\}=
\{(c[1,1,1]^\top,c[1,1,1]^\top):c\in\R\}$. The coupling forces both layers to
hold the same harmonic cycle, but cannot destroy it.

\clearpage
\section{Algorithmic Details and Pseudocode}
\label{app:algorithms}

\subsection{Incremental Update Algorithm}

\begin{algorithm}[h]
\caption{Incremental Supra-Hodge Spectral Update}
\label{alg:update}
\begin{algorithmic}[1]
\Require Current $L_p$, perturbation $\Delta L_p\succeq 0$, current $\Up$
\Ensure Updated eigenvectors $\Up'$ (warm-started)
\State $\Delta\Lambda \gets \Up^\top(\Delta L_p)\Up$ \Comment{First-order eigenvalue shifts}
\State $\Lambda' \gets \Lambda + \diag(\Delta\Lambda)$ \Comment{Updated eigenvalues}
\State Initialize Lanczos at $\Up$ on $L_p+\Delta L_p$
\State Run $t_{\max}$ Lanczos steps to refine $\Up'$
\State \Return $\Up'$, $\Lambda'$
\end{algorithmic}
\end{algorithm}

For small $\Delta L_p$ (rank-one simplex addition), this requires $O(r\cdot kN_p)$
operations versus $O(r\cdot kN_p\log(kN_p))$ for a full restart.

\clearpage
\section{Experiment-Generation Details}
\label{app:experiments}

\textbf{Random complex generation.} For each trial, generate $n$ entities with
embeddings $e_v\sim\mathcal{N}(0,I_{32})$, normalized to unit sphere.
Similarity: $\simop(u,v)=e_u^\top e_v$. Triplet score:
$\tri(u,v,w)=\min(\simop(u,v),\simop(u,w),\simop(v,w))$. Thresholds:
$\tau_1=0.6$, $\tau_2=0.5$.

\textbf{Planted cycles.} For each planted cycle, select a random unfilled
triangle $\{u,v,w\}$ in a randomly chosen layer $i$ where all three edges
$\{u,v\},\{u,w\},\{v,w\}$ are present but the 2-simplex $\{u,v,w\}$ is absent.
Record the set of planted (edge, layer) pairs as ground truth.

\textbf{Detection criterion.} Let $u_1,\dots,u_r$ denote eigenvectors of $L_1$
associated with the $r$ smallest eigenvalues, for fixed small $r$ (we use $r=5$).
For each planted cycle indicator vector, check whether at least one of these
low-energy eigenvectors has cosine similarity greater than $0.8$ with that
indicator. If so, that planted cycle is counted as detected.

\textbf{Compression baseline.} The PCA baseline uses the top-$r$ SVD of the
stacked signal matrix $[v^{(1)}\cdots v^{(k)}]\in\R^{N_p\times k}$; the
Supra-Hodge baseline uses $\Pp$.

%% ── Bibliography ──────────────────────────────────────────────────────────────
\begin{thebibliography}{99}

\bibitem{chung1997}
F.R.K. Chung,
\textit{Spectral Graph Theory},
CBMS Regional Conference Series in Mathematics, Vol.~92,
American Mathematical Society, 1997.

\bibitem{gomez2013}
S. G{\'o}mez, A. D{\'i}az-Guilera, J. G{\'o}mez-Garde\~{n}es,
C.J. P{\'e}rez-Vicente, Y. Moreno, A. Arenas,
Diffusion dynamics on multiplex networks,
\textit{Physical Review Letters}, 110(2):028701, 2013.

\bibitem{dedomenico2013}
M. De Domenico, A. Sol{\'e}-Ribalta, E. Cozzo, M. Kivel{\"a}, Y. Moreno,
M.A. Porter, S. G{\'o}mez, A. Arenas,
Mathematical formulation of multilayer networks,
\textit{Physical Review X}, 3(4):041022, 2013.

\bibitem{horak2013}
D. Horak, J. Jost,
Spectra of combinatorial Laplace operators on simplicial complexes,
\textit{Advances in Mathematics}, 244:303--336, 2013.

\bibitem{schaub2020}
M.T. Schaub, A.R. Benson, P. Horn, G. Lippner, A. Jadbabaie,
Random walks on simplicial complexes and the normalized Hodge Laplacian,
\textit{SIAM Review}, 62(2):353--391, 2020.

\bibitem{lim2020}
L.-H. Lim,
Hodge Laplacians on graphs,
\textit{SIAM Review}, 62(3):685--715, 2020.

\bibitem{yao2023}
S. Yao, J. Zhao, D. Yu, N. Du, I. Shafran, K. Narasimhan, Y. Cao,
ReAct: Synergizing reasoning and acting in language models,
\textit{International Conference on Learning Representations (ICLR)}, 2023.

\bibitem{yang2018}
Z. Yang, P. Qi, S. Zhang, Y. Bengio, W.W. Cohen,
R. Salakhutdinov, C.D. Manning,
HotpotQA: A dataset for diverse, explainable multi-hop question answering,
\textit{Proceedings of EMNLP}, 2018.

\bibitem{edge2024}
D. Edge, H. Trinh, N. Cheng, J. Bradley, A. Chao, A. Mody, S. Truitt, J. Larson,
From local to global: A graph RAG approach to query-focused summarization,
\textit{arXiv:2404.16130}, 2024.

\bibitem{qin2023}
Y. Qin, S. Hu, Y. Lin, W. Chen, N. Ding, G. Cui, Z. Zeng, Y. Zhou, X. Huang,
Z. Liu, M. Zheng, Y. Cai, C. Liu, D. Suyal, Z. Liu, M. Sun, Z. Liu,
Tool learning with foundation models,
\textit{arXiv:2304.08354}, 2023.

\bibitem{hansen2020}
J. Hansen, T. Gebhart,
Sheaf neural networks,
\textit{Advances in Neural Information Processing Systems}, 33:10437--10448, 2020.

\bibitem{bodnar2021}
C. Bodnar, F. Frasca, Y. G. Wang, N. Otter, P. Li{\'o}, M. M. Bronstein, M. Deac,
Weisfeiler and Lehman go topological: Message passing simplicial networks,
\textit{Proceedings of ICML}, 2021.

\end{thebibliography}

\end{document}